% ============================================================
% BlackboxNLP 2026 submission — working draft
% Rewritten july3 (streamlined arc); previous version: main-old.tex.bak
%
% Build: pdflatex main / latexmk -pdf main
% NOTE: plain article class for drafting. Before submission switch to the
% official ACL style (acl-style-files); BlackboxNLP uses the ACL template.
%
% Convention: the abstract is prose (Hadi's); everything else is \outline{}
% bullets that Hadi converts to prose. Claude maintains structure/refs/numbers.
% Page budgets in section headings (X.Xp) assume ACL LONG format: 8 pages of
% content + unlimited Limitations/references. Current allocation ~8.5p counted
% -> trim ~0.5p in the prose pass (Background and Pilot are the flex).
% Bullet tags: (half para) / (1 para) / (2 paras) = target length in prose.
%
% Narrative arc: phenomenon (unexplained steering variance) -> inventory
% (~600 concepts, 24 supersenses, dual contrasts) -> two-config steering ->
% regression (which concepts steer) -> pilot interpretation (parent-dominance;
% coherence != identity) -> causal probe (within-category rescue).
% ============================================================
\documentclass[11pt]{article}
\usepackage[margin=1in]{geometry}
\usepackage{enumitem}
\usepackage{xcolor}
\usepackage{amssymb}
\usepackage{natbib}
\usepackage[normalem]{ulem}  % \sout{...} available; convention: strike, don't delete

% outline bullets render gray so they're visually distinct from prose
\newenvironment{outline}{\color{gray!70!black}\begin{itemize}[leftmargin=1.2em,itemsep=1pt]}{\end{itemize}\color{black}}

\title{Not Every Concept Is a Direction:\\
Which Lexicalized Concepts Can Be Steered in LLMs?
\thanks{Working title. Alternative: ``Concreteness and Polysemy Predict the Steerability of Lexicalized Concepts'' (use if R2 confirms). \textbf{[NEW jul6] R2 did NOT confirm: polysemy is null and concreteness is negative; frequency is the positive predictor. Do not use the alternative title.}}}
\author{Hadi Asghari \and ...}
\date{Draft of \today}

\begin{document}
\maketitle

\begin{abstract} % (0.5p incl. title block)
Steering vectors are widely used to control and interpret LLMs, yet benchmarks show that some concepts steer reliably while others barely respond --- a variance that is reported but rarely explained. In this research we investigate which kinds of concepts are steerable in Llama-3.1-8B and Gemma-2-9B. We use the RFM method to extract steering vectors (or directions) for approximately 600 concepts drawn from 24 WordNet supersense categories, and validate which of them successfully steer the model. We then build a regression model explaining steerability from concept category, concreteness, polysemy, and word frequency, as well as the share of the top eigenvector during extraction ($\lambda_1$/trace). The results help distinguish concepts that behave as single directions --- ``primitive'' for the model --- from those that are composite or context-dependent, and offer practitioners an extraction-time indication of when single-vector steering can be trusted.

\begin{outline}
\item \textbf{[NEW jul6, replaces the expectation below] Main run (600 concepts, Llama) is in: word frequency is a \emph{positive} predictor ($\beta_{\mathrm{zipf}}{=}+0.40$, $p{=}.02$) and concreteness a \emph{negative} one ($\beta_{\mathrm{conc}}{=}-0.28$, $p{=}.017$); polysemy and AoA are null. The pilot expectation (concreteness+polysemy carry it, frequency null) is rejected on both counts: frequent-but-abstract concepts steer best.}
\item Candidate sentence once results are in (pilot expectation): concreteness and polysemy --- \emph{not} frequency --- predict steerability; the frequency null is itself a finding. (1 sentence)
\item Candidate sentence (pilot mechanism): when steering fails, it fails informatively --- single-vector steering delivers the concept's \emph{category prototype} (a city steers to ``a beloved place'', not to itself). (1 sentence)
\item Candidate sentence (causal probe): re-extracting with within-category contrasts tests whether such failures lie in the extraction design rather than the model. (1 sentence)
\item Candidate sentence (replication): findings replicate across two model families (Llama, Gemma). (1 sentence)
\end{outline}
\end{abstract}

\section{Introduction (1.0p)}
\begin{outline}
\item Hook: steering/RePE papers throw heterogeneous concept lists at models; some concepts steer at 90\%, others at chance. The variance is treated as noise; we argue it is signal. (1 para)
\item The implicit assumption of top-1 eigenvector / PCA-1 / difference-of-means steering: the concept is representable as \emph{one direction}, injectable with \emph{one coefficient}, at whatever layer(s) the method uses. (half para)
\item Existing evidence that not all features are one-dimensional (Engels et al.) --- but no one has asked \emph{which} concepts behave as single directions, or whether this aligns with human semantic structure. (half para)
\item Our question: is steerability systematic and predictable from the linguistic properties of the word --- and what does a steering failure actually consist of? (half para)
\item Contributions (short lead-in + three numbered items, 2--3 sentences each): (1 para)
  \begin{itemize}
  \item C1: a systematic steerability study --- $\sim$600 lexicalized concepts across 24 WordNet supersenses, annotated with psycholinguistic norms, steered under a controlled two-configuration protocol; per-concept results and code released.
  \item C2: a regression explaining steerability variance from concept category, concreteness, polysemy, and frequency, with extraction-internal spectral metrics ($\lambda_1$/trace, $\lambda_1/\lambda_2$) as candidate zero-cost predictors.
  \item C3: a mechanistic account of what failure looks like, grounded in pilot diagnostics (parent-dominance of extracted directions; coherence $\neq$ identity) and tested with one causal probe (within-category contrast rescue).
  \end{itemize}
\end{outline}

\section{Background and Related Work (1.0p)}
% Convention per bullet: cited line(s) + one closing sentence on same-vs-different.
\begin{outline}
\item \textbf{Linear representation hypothesis and steering.} Park et al.; RepE (Zou et al.); CAA (Panickssery/Rimsky et al.); ActAdd (Turner et al.); ITI (Li et al.); the RFM/universal-steering benchmark we build on (Beaglehole et al.\ 2026) --- kernel ridge + AGOP metric learning, top eigenvector as steering vector. \emph{Ours:} we adopt this family's method (RFM) unchanged, but instead of proposing a steering technique we ask \emph{which concepts} the technique works for and why. (1 para)
\item \textbf{Multi-dimensional features.} Engels et al.\ (not all features are one-dimensional; circular features); feature splitting in SAEs; conceptors / subspace steering. \emph{Ours:} same premise (one direction may not suffice), but we test it at the level of a stratified concept inventory and connect it to word properties rather than to individual features. (half para)
\item \textbf{[NEW jul6] Hierarchical concept geometry --- must cite.} Park, Choe, Jiang \& Veitch (ICLR 2025, arXiv:2406.01506): $\sim$900 WordNet concepts on Llama-3 and Gemma; categorical concepts are simplices and hierarchically related concepts occupy \emph{orthogonal} subspaces. The representational-geometry twin of this study (same inventory source, same model families). \emph{Ours:} they characterize how the hierarchy is \emph{encoded}; we test what the hierarchy predicts about \emph{intervention} --- and our parent-dominance finding (extraction returns the shared category axis) is the steering-side counterpart of their parent/child orthogonal decomposition; R4 tests whether subtracting the parent axis recovers the child. (half para)
\item \textbf{Predicting steerability.} LAP / logit-lens accessibility (Billa 2026); ``When is your LLM steerable?''; (un)reliability of steering vectors. \emph{Ours:} these predict from engineering signals; we predict from the \emph{linguistic structure} of the concept, and additionally explain what the failure mode consists of. (half para)
\item \textbf{[NEW jul6] Add cites here:} (a) ``Understanding (Un)Reliability of Steering Vectors in Language Models'' (arXiv:2505.22637) --- directional agreement between per-example activation differences and the steering vector predicts steerability; (b) its 2026 follow-up on \emph{geometric predictors} of steering unreliability (arXiv:2602.17881) --- closest neighbor to our R3; differentiate: they predict from activation geometry, we from the concept's linguistic structure plus extraction-time spectra, and we characterize the failure mode (parent delivery). (c) AxBench (arXiv:2501.17148) --- the standard large-scale steering evaluation (prompting/finetuning baselines beat SAE steering); one line positioning our judge protocol and per-concept inventory relative to it.
\item \textbf{Causal identification of concept vectors.} Function vectors (Todd et al.\ 2024); linear relation embeddings (Hernandez et al.\ 2024); ROME (Meng et al.) --- identification by \emph{intervention}. \emph{Ours:} complementary --- we identify by \emph{contrast design} and use one intervention-style probe (R4) to arbitrate when the contrast is at fault. (half para)
\item \textbf{Frequency and linear representations.} Co-occurrence thresholds for linear representations; long-tail knowledge (Kandpal et al.). \emph{Ours:} we test the frequency hypothesis directly and (per pilot) expect to reject it --- exposure alone does not buy steerability. (half para)
\item \textbf{LLM--human concept alignment and psycholinguistic norms.} Concreteness/AoA norms in NLP; Shani et al.\ 2025 (``From Tokens to Thoughts'', arXiv:2505.17117): LLMs trade compression for meaning differently than humans, measured on \emph{input-side} token embeddings vs.\ human categories. \emph{Ours:} same alignment question, but measured mid-stream and \emph{causally} --- does human semantic structure predict what can be steered, not just how embeddings cluster. (half para)
\item \textbf{Methodological caution on cosine.} You 2025 (``Semantics at an Angle'', arXiv:2504.16318): cosine similarity can mislead in anisotropic learned spaces (a shared dominant component inflates all similarities). \emph{Ours:} our geometry claims rest on \emph{differential} cosines under identical measurement (within-class 0.84 vs.\ between-class 0.35 vs.\ moods 0.46), with the between-class value reported as the anisotropy control; mid-stream RMSNorm anchors the metric. (half para)
\item Gap statement: no work connects concept-level semantic structure (concreteness, polysemy, taxonomic class) to steering outcome, nor explains what an unsteerable concept's failure consists of. (half para)
\end{outline}

\section{Defining and Sampling Concepts (1.0p)}
\begin{outline}
\item \textbf{Operational definition.} Sidestep the philosophical debate (Murphy 2002; Margolis \& Laurence 1999): a \emph{lexicalized concept} = a meaning a speech community assigned a word to, operationalized as a WordNet synset; each word mapped to its dominant (first) noun sense. (half para)
\item \textbf{Inventory.} All single-word, noun-dominant lemmas from WordNet; category = supersense (26 lexicographer files; Miller 1990; Fellbaum 1998; ``supersense'': Ciaramita \& Johnson 2003). Exclusions: \texttt{noun.Tops}; POS-ambiguous lemmas; small cells sampled exhaustively and flagged. (half para)
\item \textbf{Sample ($\approx$600, flat).} 25/supersense across $\sim$24 usable supersenses; within each supersense, stratify over frequency so category and frequency are not confounded. Six supersenses are pre-specified (before extraction) as the \emph{featured} cells for within-category analyses: \texttt{location}, \texttt{person}, \texttt{feeling} (pilot continuity) and \texttt{artifact}, \texttt{cognition}, \texttt{food} (concreteness extremes). At $n{=}25$ every cell is adequately powered for class means and absorbs attrition (prompted-ceiling/judge exclusions). (1 para)
\item \textbf{Predictors} (one table): Zipf frequency (wordfreq); concreteness (Brysbaert et al.\ 2014); age of acquisition (Kuperman et al.\ 2012); polysemy (\#WordNet senses); taxonomic depth (covariate); supersense; log category size (covariate). (half para + table)
\item \textbf{Contrast design.} Positives wrap a statement in a persona frame; the frame must be felicitous across \emph{all} supersenses (e.g.\ ``You are fascinated by \{concept\}''), else frame felicity confounds category. \emph{Dual negatives}: (i) bare statement (benchmark-comparable); (ii) same frame with a \emph{different} concept --- free, since those prompts are other concepts' positives --- which cancels the frame axis, and with same-supersense negatives also the category parent. Rationale: a contrast in which frame, category, and instance always co-occur cannot statistically separate them (see Pilot). (1 para)
\item \textbf{Pilot/confirmatory split.} Pilot = 90 concepts (30 each: places, moods, personalities) from the released benchmark lists, used for all exploration; the WordNet sample is confirmatory, untouched during exploration. (half para)
\end{outline}

\section{Method (1.0p)}
\begin{outline}
\item \textbf{Direction extraction (RFM) --- and why RFM.} Per layer: kernel ridge regression on concept-vs-negative prompts; AGOP $M = \mathbb{E}[\nabla f \nabla f^{\top}]$; steering vector = top eigenvector of $M$. Chosen among RepE methods because (a) it is the method of the benchmark we build on (comparability); (b) unlike difference-of-means or PCA-1 it returns the \emph{full discriminative spectrum} with eigenvalues --- the spectral metrics below come for free, and the method is self-diagnosing (the pilot's confound was discovered through them); (c) its steering performance matches or beats the simpler extractors on the released benchmark. (1 para)
\item \textbf{Spectral metrics.} $\lambda_1/\mathrm{tr}(M)$ and $\lambda_1/\lambda_2$; per-concept score = band average over decoder depth 4--24 (per-layer profile retained). These measure the \emph{coherence} of the extracted signature, not concept identity (see Pilot). (half para)
\item \textbf{Fixed hyperparameters --- a methodological requirement.} Kernel bandwidth inflates $\lambda_1$/trace mechanically (large bandwidth $\Rightarrow$ near-linear fit $\Rightarrow$ rank-1 AGOP); grid search over near-tied validation correlations injects artifact variance into the dependent measure. All extractions at bw=1, norm=True, reg=$10^{-3}$; vector and metrics read at a fixed (converged) iteration. (half para)
\item \textbf{Two-config steering protocol.} Each concept steered under (i) \emph{all layers}, benchmark defaults (unit vectors, released coefficient grid --- comparability); (ii) \emph{one mid layer} ($\approx$40\% depth; vector scaled by the concept's per-layer projection gap, cf.\ CAA / ITI; effective single-site dose $\approx$1$\times$ local residual-stream norm). $\Delta$(all $-$ single) measures how \emph{distributed} the concept's steerable content is. 3 category-neutral prompts (listed verbatim); greedy decoding. (1 para)
\item \textbf{Outcome.} Primary: fraction of (coef $\times$ prompt) cells judged as expressing the concept; explicitly \emph{not} max-over-coefficients. Secondary: \emph{granularity reached} (category / region / instance) --- the binary criterion flattens near-misses (pilot: Edinburgh$\to$``Scottish Highlands'' scores 0). (1 para)
\item \textbf{[NEW jul6] Graded judge (V8).} The 600-run generations show systematic sub-threshold movement the binary judge cannot credit (ragtime $\to$ lyre/melody imagery; candlepower $\to$ binoculars/glint): a graded rubric (1.0 named / 0.75 immediate field / 0.5 category level / 0.25 trace / 0 unrelated) re-judges the same generations and yields a continuous steerability outcome for the regression (thin-tail fix: only 119/600 binary positives).
\item \textbf{Judge.} gpt-oss:20b (local, on the experiment GPU); symmetric judge question across supersenses. Validated against Claude Sonnet~5 on a 360-item set: 87--91\% agreement, $\kappa = .72$--$.78$; a stricter candidate (Haiku 4.5) was rejected for one-directional divergence from the two-judge consensus. Confirmatory-study additions: re-validation on $\sim$200 WordNet-distribution items + human spot-check of disagreements (Appendix~E). (half para)
\item \textbf{Controls.} Prompted ceiling (can the model express the concept when explicitly asked?) --- steerability conditional on ceiling; pipeline sanity-checked on known-good concepts per model. (half para)
\item \textbf{Models.} Full protocol on Llama-3.1-8B-Instruct; headline analyses replicated on Gemma-2-9B-it. (half para)
\end{outline}

\section{Pilot Diagnostics and Protocol Selection (1.0p)}
% All pilot experiments: N = 30x3 (places/moods/personalities), one 90-concept sample.
% Behavioral arms re-run at N=90 pending; interim n=15 shakedown numbers are provisional.
\begin{outline}
\item \textbf{Hyper-parameter selection (summary).} Validation $r$ saturates ($\sim$0.97--0.99) across bandwidths, statement counts, and iterations $\Rightarrow$ selection by validation is third-decimal noise; large bandwidth forces rank-1 AGOP by construction; converged readout cuts metric instability 3--4$\times$. Decision rule: fixed modal setting, converged readout. Details in Appendix~A. (half para)
\item \textbf{Finding 1: parent-dominance.} Different same-category concepts' extracted top-1 directions are nearly identical (places: cosine $0.97 \to 0.84 \to 0.72$ across depth; \emph{all} top-3 components $\geq$0.90 shared). Single-layer dose-response saturates at the category \emph{prototype} (every city $\to$ ``quiet beach at sunset'') with region-level near-misses (Edinburgh$\to$Dublin pub); giving secondary components injection mass adds a generic roleplay axis, not instance content. Three legs: geometry, dose-response, component ablation. Interpretation: a contrast in which frame, parent, and instance always co-occur is statistically one axis --- what steering delivers is mostly the category. (1 para + figure \texttt{parent\_share\_by\_depth.png})
\item \textbf{Finding 2: coherence $\neq$ identity.} Places show the \emph{highest} band $\lambda_1$/trace yet fail single-layer steering; moods (lower) succeed. Spectral metrics certify that a \emph{coherent} (context-invariant, rank-1) direction was extracted --- not that it is the judged content. Consequence: their regression coefficients are expected to be conditional on contrast type; tested under both contrasts. (1 para)
\item \textbf{Finding 3: attribute vs.\ instance.} Attribute concepts (moods) steer from one mid layer; instance identities steer only under all-layer injection, reaching region granularity --- consistent with multi-stage assembly (per-layer residue accumulation + late lexical tipping + autoregressive bootstrap once the name is emitted). Motivates the two-config protocol and $\Delta$. (1 para)
\item \textbf{Protocol justifications} (one line each; evidence in Appendices~C--D): single mid layer at $\approx$40\% depth --- wide behavioral window there; 60\% yields token-bias then repetition collapse. No top-$k$ arms --- the top-3 subspace under the released contrast is category-shared. Coefficient $\approx$1$\times$ local residual norm via per-layer magn scaling. Granularity as secondary outcome --- the binary judge flattens near-misses. (1 para)
\end{outline}

\section{Results (1.5p)}
\begin{outline}
\item \textbf{R1: Steerability varies systematically.} Steering success varies widely across concepts and clusters by supersense; distribution figure per supersense, per config. (half para + figure)
\item \textbf{[NEW jul6] R1 numbers (Llama, all-layer config, v2b frame, judge = gpt-oss).} 81/600 concepts steer under the subject-invite evaluation (V1), 95/600 under the neutral evaluation (V6), 119/600 (19.8\%) under either; V1--V6 agreement $\phi = 0.59$ ($\kappa=0.59$) --- substantial but partly evaluation-relative, motivating the union outcome and the graded judge (V8). Per-supersense success spans 2--9 of 25 (no ceiling/floor): top \texttt{noun.feeling} (9), \texttt{noun.communication} (8), \texttt{noun.food} (7); bottom \texttt{noun.body} (2), \texttt{noun.plant} (4), \texttt{noun.artifact} (4). Figures: \texttt{plots/steered\_by\_supersense\_600.png}, \texttt{plots/steered\_boxplots\_600.png}; analysis in \texttt{nb\_analysis600.ipynb}.
\item \textbf{R2: Psycholinguistic predictors.} Regression of steerability on concreteness, polysemy, AoA, Zipf, taxonomic depth (+ category size). Pilot expectation: concreteness and polysemy carry the signal; frequency $\approx$ 0 --- exposure alone does not buy steerability. (1 para + table)
\item \textbf{[NEW jul6] R2 outcome (Table~\ref{tab:reg600}): the pilot expectation is rejected on both counts.} Zipf frequency predicts steerability \emph{positively} ($\beta{=}+0.40$, $p{=}.02$) and concreteness \emph{negatively} ($\beta{=}-0.28$, $p{=}.017$, robust across specifications); polysemy and AoA are null. Frequent-but-abstract concepts steer best. Collinearity is benign among the lexical predictors (conc--zipf $r{=}.09$; all lexical VIF $<2.7$).
\item \textbf{R3: Spectral metrics as candidate predictors.} $\lambda_1$/trace and $\lambda_1/\lambda_2$ within the full regression; validity compared across the two contrast types (expected conditional --- Pilot Finding 2); controlling for prompted ceiling and val-$r$ (saturated $\Rightarrow$ cannot explain variance). Baselines: frequency alone, category alone. (1 para)
\item \textbf{[NEW jul6] R3 outcome: the robust spectral predictor is the AGOP \emph{participation ratio}, not $\lambda_1$/trace.} PR is negative both alone ($\beta{=}-0.27$, $p{=}.023$) and in the full model ($\beta{=}-0.72$, $p{=}.007$): concepts whose discriminative spectrum concentrates on fewer directions steer better. $\lambda_1$/trace has \emph{no} independent effect (n.s.\ alone, sign-unstable in the trio; VIF $=12$ within the spectral set) --- quantitative confirmation of Pilot Finding 2 (coherence $\neq$ identity) at $n{=}600$: top-1 dominance per se does not buy steering. Report the one-spectral-at-a-time ladder alongside the full model (Table~\ref{tab:reg600}).
\item \textbf{[NEW jul6] Category is mediated by measurable word properties.} Supersense dummies do not pay for themselves: AIC 626.7 (supersense-only) vs.\ 599.7 (null) vs.\ 576.7 (lexical+spectral, best); adding supersense to the best model \emph{worsens} AIC (597.9). The category signal visible in R1 is carried by the continuous predictors --- a stronger claim than ``categories differ.''
\item \textbf{R4: The rescue probe (causal).} For instance-like concepts that fail: re-extract with within-category contrasts (``loves Edinburgh'' vs.\ ``loves \{other city\}''; frame and parent cancel) and steer at one mid layer. Success $\Rightarrow$ the failure was extraction addressing (contrast design), not model capability; failure $\Rightarrow$ instance content is not single-site injectable (assembly account). (1 para)
\item \textbf{TODO [jul6]: R4 has not been run yet --- it is the highest-priority remaining experiment.} Everything else in this section now has data (R1--R3, R5 at $n{=}600$; R6 pilot + Gemma replication in progress); R4 is the paper's one \emph{causal} claim and is cheap (re-extract $\sim$5--25 concepts at $\sim$3\,s each with within-category negatives + one single-layer eval; the zero-reextraction averaging variant in Appendix~C costs nothing). Run immediately after the Gemma replication.
\item \textbf{R5: Distributedness.} $\Delta$(all-layer $-$ single-layer) per concept, reported descriptively by supersense and correlated with the R2 predictors. (half para)
\item \textbf{R6: Generality.} Llama vs.\ Gemma consistency: per-concept correlation of steerability and of the spectral metrics across models. (half para)
\item \textbf{[NEW jul6] R6 pilot evidence: behavior replicates, spectra do not.} On the 90-concept pilot (orig frames, V1) Gemma-2-9B matches Llama class-by-class (moods 30/30 both; personalities 23 vs.\ 24; places 0 vs.\ 2), with coefficients calibrated to Gemma's $\sim$50$\times$ larger residual norms (usable window 12--21). But the $\lambda_1$/trace class \emph{ordering inverts} across models (Llama: places $>$ pers $>$ moods; Gemma: moods $\approx$ pers $>$ places) --- spectral coherence is partly model-geometry, not concept-intrinsic, so cross-model spectral comparisons need within-model standardization. Gemma 600-concept replication in progress.
\end{outline}

% [NEW jul6] regression table for R2/R3 (numbers from nb_analysis600.ipynb; logit,
% z-scored predictors, outcome = steered under V1 or V6, n=600)
\begin{table}[t]
\centering
\small
\begin{tabular}{lrrr}
\hline
 & lexical & lexical & lexical \\
 & only & $+$PR & $+$spectral \\
\hline
concreteness   & $-0.20^{\dagger}$ & $-0.28^{*}$  & $-0.45^{**}$ \\
Zipf frequency & $+0.40^{*}$       & $+0.40^{*}$  & $+0.33^{\dagger}$ \\
AoA            & $-0.15$           & $-0.08$      & $-0.17$ \\
polysemy       & $-0.03$           & $-0.02$      & $+0.08$ \\
n tokens       & $+0.08$           & $+0.07$      & $+0.03$ \\
$\lambda_1$/trace   &              &              & $-0.18$ \\
$\lambda_1/\lambda_2$ &            &              & $-0.44^{\dagger}$ \\
participation ratio &              & $-0.27^{*}$  & $-0.72^{**}$ \\
\hline
AIC            & 591.3             & 588.0        & 576.7 \\
\hline
\end{tabular}
\caption{\textbf{[NEW jul6]} Logistic regression of steerability (V1 $\cup$ V6) on
z-scored predictors, $n{=}600$ (Llama-3.1-8B, all-layer, v2b frame). Reference AICs:
null 599.7; supersense-only 626.7; supersense $+$ all predictors 597.9.
$^{\dagger}p{<}.1$, $^{*}p{<}.05$, $^{**}p{<}.01$. Spectral metrics are
inter-correlated (PR--$\lambda_1$/trace $r{=}-.85$); the PR-only column is the
conservative specification.}
\label{tab:reg600}
\end{table}

\section{Discussion (1.0p)}
\begin{outline}
\item What makes a concept steerable: attributes (context-general colorings) admit a single injectable direction; instance identities require assembled, multi-stage content --- steering failure is usually \emph{delivery of the parent category}, not silence. (1 para)
\item What $\lambda_1$/trace is: an \emph{intervention-match score} --- how well the concept's natural activation signature fits the one move single-vector steering can make (one constant vector, all contexts). Mechanism for R2: concrete concepts express context-invariantly; abstract ones context-dependently. (half para)
\item \textbf{[NEW jul6] Revision required by R2/R3: the intervention-match story must move from $\lambda_1$/trace to the participation ratio, and the concreteness mechanism reverses.} Empirically it is spectrum \emph{concentration} (low PR), not top-1 dominance, that predicts steering, and \emph{abstract} (not concrete) concepts steer better --- consistent with the pilot's parent-dominance reading: for concrete instance-like concepts the dominant direction is largely the shared category axis, so its coherence does not deliver the instance. (half para)
\item Status of extracted vectors: the residual stream has no privileged basis, so extracted directions are \emph{experiment-relative}; within-model relational claims are valid, coordinate-level and raw cross-model comparisons are not. Anchored spaces (embedding/unembedding, MLP neurons) are where mechanism pins a frame. (half para)
\item Two routes to a vector's identity: \emph{experiment design} (contrasts that vary exactly one thing) vs.\ \emph{causal mediation} (function vectors, ROME). Pipeline: propose by contrast, accept by causal test --- R4 instantiates it. (half para)
\item Practical implications for RePE: report spectral concentration and contrast design alongside steering benchmarks; stratified concept inventories; granularity-graded judging over binary. (half para)
\end{outline}

\subsection*{Future work: abstraction-depth as a second axis (0.5p)}
\begin{outline}
\item Idea: besides \emph{how many} directions compose a concept, \emph{where} (which layer) it first becomes steerable --- hypothesis: hypernyms steer shallower than hyponyms (cf.\ pilot: moods at depth $\sim$9, places at $\sim$17+). (half para)
\item Pilot lesson: this is a \emph{write} property --- extraction-geometry proxies fail the steering ground truth; measurement requires single-layer steering sweeps. Test on the WordNet hierarchy (taxonomic depth); independence from the compositional axis must be assessed \emph{within} supersense. (half para)
\item Related vignette: compositional concepts --- is direction(``strict father'') $\approx \alpha\,$direction(``strict'') $+ \beta\,$direction(``father'')? (half para)
\end{outline}

\subsubsection*{[NEW jul6] Scoped design: first-steerable-depth sweep on the featured cells}
\begin{outline}
\item \textbf{Design.} Not $600 \times$ many layers: the six \emph{featured} supersenses ($150$ concepts; \texttt{location}, \texttt{person}, \texttt{feeling}, \texttt{artifact}, \texttt{cognition}, \texttt{food}) $\times$ $\sim$5 single-layer sites spanning depth (e.g.\ 9 / 13 / 19 / 24 / 28), magn-calibrated, generation$+$judging only --- the per-layer directions are already extracted, so each site is one eval pass ($\sim$1.2\,h local); one overnight total. (half para)
\item \textbf{Payoff.} Operationalizes the abstraction-depth axis with WordNet taxonomic depth as ground truth: \emph{where} in depth does a concept first become injectable, and does the hierarchy predict it (hypernyms shallower than hyponyms)? Preliminary $n{=}600$ evidence at two sites is consistent: depth-13 retention is ordered by abstractness (cognition 5/6, feeling 4/9, communication 4/8 survive; substance/process/artifact/body retain 0). Independence from the compositional-rank axis assessed \emph{within} supersense (the two are confounded between classes). (half para)
\item \textbf{Decision rule.} Run after R4. Present as a single figure (retention-by-depth per supersense); if the depth profiles are noisy or degenerate, the section stays Future Work and the figure moves to the appendix. (half para)
\end{outline}

\section{Limitations (0.5p; uncounted in ACL format)}
\begin{outline}
\item Single extraction method (RFM); $\lambda_1$/trace is AGOP-specific --- PCA explained-variance as cross-check (appendix if time). (half para)
\item Steering conclusions are relative to the additive intervention class; clamp/projection interventions may behave differently (targets saved, untested). (half para)
\item LLM judge (validated but imperfect); English only; dominant-sense assumption; norms are human ratings. (half para)
\item Correlational core with one causal probe (R4); extraction-design effects and model limits not fully separable outside that design. (half para)
\end{outline}

% \section*{Ethics Statement}  % likely minimal
% \section*{Acknowledgments}

\bibliographystyle{acl_natbib} % needs acl-style-files; plainnat while drafting
% \bibliography{refs}

% ============================================================
% Experiment-track appendices (working log of what we tried).
% Rewritten july3 (streamlined); previous version: appendix_experiments-old.tex.bak
% Input into main.tex before \end{document}:  % ============================================================
% Experiment-track appendices (working log of what we tried).
% Rewritten july3 (streamlined); previous version: appendix_experiments-old.tex.bak
% Input into main.tex before \end{document}:  % ============================================================
% Experiment-track appendices (working log of what we tried).
% Rewritten july3 (streamlined); previous version: appendix_experiments-old.tex.bak
% Input into main.tex before \end{document}:  \input{appendix_experiments}
% Relies on main.tex preamble (outline env, ulem, amssymb). Do NOT add a preamble here.
%
% Pilot standard: N = 30x3 on the single 90-concept sample (places/moods/personalities).
% A few early bullets cite interim n=15/class shakedown numbers (flagged in place); the
% depth sweep (C) supersedes the two-site depth conclusions at n=90.
%
% A--F support the main arc. Z1--Z3 are PARKED (record / follow-up; submission-optional).
% Tags:  [+] worked / kept   [--] didn't work / dropped   [?] mixed / open
% Image + notebook files named in \texttt{}; figures intentionally not embedded.
% ============================================================
\appendix

\section{Appendix A: Extraction and metric setup}
\begin{outline}
\item [+] Validation $r$ (probe accuracy) saturates ($\sim$0.99) across bandwidths, statement counts, and RFM iterations $\Rightarrow$ any selection \emph{by} val-$r$ is third-decimal noise; it is also a poor steerability predictor. Boxplots: \texttt{box\_rsq.png}.
\item [+] $\lambda_1$/trace at bandwidth $L{=}1$ discriminates across concepts; $L{=}10$ collapses it toward 1 (rank-1 AGOP by construction). \texttt{box\_l1trace.png}, \texttt{box\_l1l2.png}.
\item [+] The released $\lambda_1$/trace (read at the best-$r$ iteration) is unstable --- it samples a value climbing $0.0002 \to 0.17 \to 0.46$ across iterations at a noise-chosen point; reading at convergence cuts the statement-count swing 3--4$\times$. Both vector and metrics read at a fixed (converged) iteration.
\item [+] Saved per-layer extraction stats: top-3 eigenvalues (\texttt{evals}), natural concept magnitude (\texttt{magn} = raw projection gap), clamp targets (\texttt{tpos}, \texttt{tneg}), mean residual norm (\texttt{hnorm}) --- in the rfmstats pickles.
\item [?] Per-concept ``best layer'' by argmax is unstable: val-$r$ best is noise-driven; $\lambda_1$/trace and $\lambda_1/\lambda_2$ are bimodal with a degenerate spike at the first block (exclude early blocks). \texttt{besthist\_*.png}.
\item [?] Hyper-parameter sweep factors (probe set of 24 concepts, one per supersense): bandwidth $\times$ normalization $\times$ contrastive construction (now incl.\ within-category negatives) $\times$ iteration readout; criteria: stability, discrimination, semantic signal (Spearman vs.\ concreteness/polysemy), steering non-degradation. Artifacts: \texttt{hpsweep\_phase2\_*}.
\end{outline}

\section{Appendix B: Predicting steerability (pilot 3 classes + 600-concept main run)}
\begin{outline}
\item [+] Best single predictor of steering success is $\lambda_1$/trace at a mid layer (depth~9, grouped CV-AUC $\approx$0.86); $\lambda_1/\lambda_2$ similar; val-$r$ erratic (saturated). Logistic regressions in \texttt{nb\_regxp1.ipynb}; \texttt{L9\_metrics\_by\_steerability.png}.
\item [--] Multivariate regression on \emph{all} layers' metrics is uninformative --- severe multicollinearity (the residual stream is cumulative). Univariate per-layer CV-AUC is the clean read.
\item [+] \textbf{Sign flips with depth}: $\lambda_1$/trace predicts steerability \emph{positively} early (depth~9) but \emph{negatively} late (depth~29) --- late-layer concentration is lexical collapse, not concept concentration. \texttt{l1l2\_depth29\_dist.png}.
\item [?] Concept class is a confounder/mediator: most steerability variance is between-class; within-class signal is weak. Do not ``control for class'' if class is on the causal path (concreteness $\to$ signature coherence $\to$ steerability).
\item [?] \textbf{The inversion}: places have the \emph{highest} band-averaged $\lambda_1$/trace yet fail single-layer steering, while moods (lowest) succeed --- the metric certifies \emph{coherence} of the extracted blend, not its \emph{identity}. Predictor validity is contrast-dependent (main text, Pilot Finding~2).
\item [?] Participation ratio $\approx$2--3 for all concepts at all layers --- the discriminative subspace at one layer is effectively rank 2--3, bounding what top-$k$ can recover. \texttt{PR\_depth9\_hist.png}.
\item [+] Error analysis at depth~9: misclassifications symmetric, in a narrow overlap band; characterizes the $\sim$0.86 ceiling.
\item [+] \textbf{v2b-instrument regressions (Table~\ref{tab:reg600}; superseded as main numbers by the v4 runs).} Under the dual v2b frame: zipf $+0.40^{*}$, conc $-0.28^{*}$; AoA/polysemy null; PR the only working spectral metric ($-0.27^{*}$ alone; $\lambda_1$/trace null and sign-flipped --- later diagnosed as frame-axis eigenvector capture, Appendix~F); supersense dummies fail even against the null (626.7 vs.\ 599.7). Robustness checks that carried over: lexical collinearity benign (conc--zipf $r{=}.09$, VIF $<2.7$); Bayesian mixed logistic with supersense random intercept reproduces all signs; spectral trio too inter-correlated for joint reporting (VIF up to 12). \texttt{nb\_analysis600.ipynb}.
\item [+] \textbf{Outcome-robustness under the graded judge (V8; Fig.~\ref{fig:v8scatter}).} Re-running with the graded score as outcome: concreteness ($-$) and spectral concentration survive; the Zipf effect shrinks to marginal --- part of the binary-judge frequency effect is plausibly \emph{nameability} (frequent words surface verbatim and earn binary credit). Main text keeps the binary V1$\cup$V6 outcome (closest to the RFM benchmark's evaluation model). $r_{pb}(\mathrm{V6},\mathrm{V8}) = 0.81$.
\item [+] \textbf{The pilot--main-run tension on $\lambda_1$/trace, RESOLVED:} its v2b null was an instrument artifact (frame-axis eigenvector capture for weak concepts; axis-alignment diagnostic in Appendix~F); under the constant-prefix v4 frame it is significant alone ($+0.55^{**}$) with the pilot's sign (pilot: depth-9 CV-AUC $\approx$0.86, Fig.~\ref{fig:l9metrics}).
\item [+] \textbf{v4-instrument regressions (the main numbers; Table 2 in main text).} Llama: all four lexical predictors significant; each spectral metric alone beats the entire lexical block (AICs 775--787 vs.\ 799; null 816); best model lex$+$PR (763.6). Gemma replicates all lexical signs and PR at its own informative depth (60--70\%; null at Llama's 40\% convention) --- same law, model-specific depth. Univariate honesty note: polysemy and concreteness do \emph{not} separate in raw boxplots (MW $p{>}.6$; effects are conditional); zipf/AoA/PR separate cleanly in both models (main-text Figure 3). Spectral collinearity under v4 is extreme (PR--frac1 $r{\le}-.94$): one-at-a-time only.
\end{outline}

% v2b regression table (moved here from main text jul10)
\begin{table}[t]
\centering
\small
\begin{tabular}{lrrr}
\hline
 & lexical & lexical & lexical \\
 & only & $+$PR & $+$spectral \\
\hline
concreteness   & $-0.20^{\dagger}$ & $-0.28^{*}$  & $-0.45^{**}$ \\
Zipf frequency & $+0.40^{*}$       & $+0.40^{*}$  & $+0.33^{\dagger}$ \\
AoA            & $-0.15$           & $-0.08$      & $-0.17$ \\
polysemy       & $-0.03$           & $-0.02$      & $+0.08$ \\
n tokens       & $+0.08$           & $+0.07$      & $+0.03$ \\
$\lambda_1$/trace   &              &              & $-0.18$ \\
$\lambda_1/\lambda_2$ &            &              & $-0.44^{\dagger}$ \\
participation ratio &              & $-0.27^{*}$  & $-0.72^{**}$ \\
\hline
AIC            & 591.3             & 588.0        & 576.7 \\
\hline
\end{tabular}
\caption{\textbf{v2b-instrument regression (instrument-comparison arm).}
Logistic regression of steerability (V1 $\cup$ V6), $n{=}600$ (Llama, all-layer,
v2b frame). Reference AICs: null 599.7; supersense-only 626.7. The full-model
spectral coefficients are collinearity-distorted (PR--$\lambda_1$/trace
$r{=}-.85$); the PR-only column is the conservative specification.
$^{\dagger}p{<}.1$, $^{*}p{<}.05$, $^{**}p{<}.01$.}
\label{tab:reg600}
\end{table}

\begin{figure}[t]
\centering
\includegraphics[width=\columnwidth]{L9_metrics_by_steerability.png}
\caption{Appendix B: pilot spectral metrics at depth 9 by steering outcome
(grouped CV-AUC $\approx$0.86).}
\label{fig:l9metrics}
\end{figure}

\begin{figure}[t]
\centering
\includegraphics[width=0.8\columnwidth]{v6_vs_v8_scatter.png}
\caption{Appendix B: binary (V6) vs.\ graded (V8) judge outcomes on the same
600-run generations ($r_{pb} = 0.81$); 59 V6-failures reach V8 $\geq 0.5$.}
\label{fig:v8scatter}
\end{figure}

\section{Appendix C: Steering interventions (what moved the needle)}
\begin{outline}
\item [--] Single-layer steering at depth 9: moods steer, experts middling, places $\approx$0 --- places steer to a \emph{generic} place attractor, not the named city (judge scores correct, not an artifact).
\item [+] Two mid layers (9+17) was the best targeted multi-site result ($\sim$48/35 of 90); 9+26 worse (late layer contributes generic content); late-only steering fails outright.
\item [+] \textbf{Cross-layer directions are near-orthogonal} (cosine $\approx$0.07): the concept is re-encoded per layer, so two layers are complementary slices while extra components at one layer are not.
\item [--] Single-site depth-19: \emph{no behavioral window} --- default $\to$ lexical word-lists (bitter $\to$ ``bitter, cold, dry, sarcastic'') $\to$ repetition collapse within 0.5--2$\times$ local stream norm; moods 10/15 (partly lexical credit), places/personalities $\approx$0. Late injection biases output \emph{tokens}, not behavior.
\item [+] Single-site depth-13: \emph{wide} behavioral window (embodied persona, clean saturating dose-response, no collapse): moods 15/15; places 0/15, saturating at the category prototype (``quiet beach at sunset'' for every city) with region-level near-misses (Edinburgh$\to$Dublin pub; genomicist$\to$neuroscientist). Coefficient is not the limiting factor; \emph{vector content} is.
\item [--] Top-$k$ at one layer, $k{=}3$: with eigenvalue weights $\approx k{=}1$; with equal weights \emph{worse} (moods 15$\to$10) --- components 2--3 are a generic roleplay/self-narration axis (``my name is Luna''), 0.92/0.90 shared across cities. \textbf{Parent-dominance}: all top-3 components $\geq$0.90 shared across same-category concepts at every depth (places cosine $0.97\to0.84\to0.72$; moods individuate to 0.46) --- the contrast, in which frame, parent, and instance always co-occur, is statistically one axis. \texttt{plots/parent\_share\_by\_depth.png}.
\item [+] All-layer steering (benchmark defaults): 26/45 v1; places reach \emph{region} granularity (Edinburgh$\to$Scottish Highlands with Scots dialect; Austin named exactly) --- per-layer instance-residue accumulation + late lexical tipping + autoregressive bootstrap (once the name is emitted, context sustains it).
\item [?] Sample note: the depth-13/19, $k{=}3$, and all-layer bullets above report the interim $n{=}15$/class shakedown under the \emph{orig} frames; the depth sweep below re-measures single-site behavior at $n{=}90$ (v4 directions) across ten depths and supersedes the two-site depth conclusions. Qualitative patterns (prototype delivery, token-bias late) held.
\item [?] \textbf{Planned probe --- within-category contrast extraction} (= R4, full variant, not yet run): positives ``loves Edinburgh'', negatives ``loves \{other city\}'' (same frame and statements both sides) --- frame and parent cancel; the only label-varying signal \emph{is} the city. Steer single-layer (magn-calibrated, in-window) on $\sim$5 cities. Cities named $\Rightarrow$ failure was extraction \emph{addressing}; still generic $\Rightarrow$ instance content is not single-site injectable (assembly account). \textbf{The zero-reextraction averaging variant ($v_{\mathrm{spec}}$) has been run --- full write-up moved to Appendix~Z3.}
\item [+] \textbf{[NEW jul10] Depth sweep (pilot 90, v4 directions, single-layer magn, depths 10--19).} Totals 31/33/23/38/64/71/54/61/42/34 (d10$\to$d19; V1$\cup$V6). Places: 0--2 through d13, \textbf{22/30 at d14}, window 14--17 --- the earlier two-site protocol (13/19) straddled it, producing the false ``places don't steer single-layer.'' Moods: onset $\le$10, plateau 25--30 from d13 through d19. Personalities: no single-layer home (peak 20, collapse to 1 at d19). Peak single-site d15 = 71/90 vs.\ all-layer 77. Per-layer read metrics (val-$R^2$, frac1, PR, magn/hnorm) rise into the window but stay flat through the late decline ($|r|{\approx}.5$--$.6$ with totals): the window is not recoverable read-side. \texttt{plots/depth\_sweep\_fv4.png}, RUNLOG 021.
\item [+] \textbf{[NEW jul10] Single-layer arm of the v4 600-run} (depth 13, magn): 90/600 (15.0\%) vs.\ all-layer 249; ratio 2.8$\times$ (v2b: 3.5$\times$); 11 single-layer-only rescues (abstract/attribute skew). NB depth 13 sits below the entity window found by the sweep --- the confirmatory single-site arm should move to depth 15. RUNLOG 020.
\end{outline}

\section{Appendix D: Coefficient calibration (systematic magnitude)}
\begin{outline}
\item [--] Fixed/hand-swept coefficients are brittle: the working value shifts $\sim$10$\times$ with depth and with the number of steered layers.
\item [+] \texttt{magn} (natural projection gap, saved per layer) grows $\sim$exponentially with depth; scaling each layer's vector by \texttt{magn} makes the coefficient a dimensionless multiplier (cf.\ CAA mean-difference magnitude; ITI $\sigma$-scaling). Fig.~\ref{fig:magnbylayer}.
\item [+] $\sqrt{N}$ count-normalization (cross-layer directions $\approx$orthogonal $\Rightarrow$ quadrature) makes the multiplier roughly transfer across layer counts; usable band centers near $\sim$1.
\item [+] Single-site calibration: effective single-layer steering needs $\approx$1.0$\times$ the local residual norm (\texttt{hnorm}); \texttt{magn}/\texttt{hnorm} $\approx$ 0.4--0.5 mid-stream $\Rightarrow$ multiplier $\approx$2--3. Empirically \texttt{magn} $\approx$ norm-relative scaling ($r{=}0.90$ with \texttt{hnorm}), so the two schemes are redundant.
\item [?] Clamp / project-to-target (set the projection to the concept's natural positive-class value; \texttt{tpos}/\texttt{tneg} saved) is the one structurally different alternative --- bounded, token-dependent, self-allocating across layers; untested.
\item [?] Process lesson: three config-state incidents (stale judge CSVs, judge-cache pickles, stale coefficient list) share one root cause --- hand-set knobs without provenance. Fix: auto-derive coefficient grids from saved per-layer stats; encode a run tag in artifact filenames.
\item [+] \textbf{[NEW jul10] Gemma dose-grid shift by dose-response audit.} Reading which coefficient produced each v2b-600 success: successes concentrate monotonically toward the top of the [12,15,18,21] grid (concept surfaces at 12 for only 1--3 concepts; \emph{uniquely} at 21 for 12--13), and degeneration (repetition $>$25\%) only begins at 21 (2.5--4\% of outputs) --- the grid truncated the dose-response on its rising side. v4 600-run grid shifted to [14,17,20,23]; best-over-coefficients scoring makes the added top dose risk-free.
\item [+] \textbf{[NEW jul6] Gemma-2-9B cross-model calibration.} Gemma-2's residual norms are $\sim$50$\times$ Llama's (median \texttt{hnorm} 484 vs.\ 9.4; $\sqrt{d}$ embedding scaling), yet the usable all-layer coefficient window is 12--21, \emph{not} $50\times$ Llama's 0.5--0.8: injection compounds over 41 layers and generation degenerates into repetition above $\sim$25. Grid [12, 15, 18, 21] reproduces the Llama pilot class-by-class (moods 30/30, personalities 23/30, places 0/30) --- evidence that norm-relative reasoning brackets but does not pin the coefficient; a short empirical sweep remains necessary per model family.
\end{outline}

\begin{figure}[t]
\centering
\includegraphics[width=0.8\columnwidth]{magn_by_layer.png}
\caption{Appendix D: the per-layer natural projection gap (\texttt{magn}) grows
$\sim$exponentially with depth, motivating magn-scaled dosing.}
\label{fig:magnbylayer}
\end{figure}

\section{Appendix E: Judge selection and validation}
\begin{outline}
\item [+] Judge shootout on a fixed 360-item set (the pilot all-layer run's generations, v1+v4): local gpt-oss:20b 26/45 (v1), pod gpt-oss:20b 27/45, Claude Sonnet~5 26/45 --- with \emph{identical} per-class breakdowns for local-oss and Sonnet (10/15/1). Judge is reproducible across machines and agrees with an independent frontier model.
\item [--] Claude Haiku 4.5 is a systematically \emph{strict outlier} (13/45 v1; moods 15$\to$5): of 49 disagreements with Sonnet, 48 are one-directional (Haiku under-credits). Cause: it enforces fluency/coherence as a hard criterion, failing fragmented-but-clearly-on-concept outputs. Rejected --- and a useful demonstration that judge choice can move absolute rates $\sim$2$\times$ while the class structure survives.
\item [+] Primary-judge validation numbers (gpt-oss vs.\ Sonnet, $n{=}360$): agreement 86.7\%/90.6\% (v1/v4), Cohen's $\kappa = 0.72/0.78$ --- substantial agreement. Artifacts: \texttt{oss\_judgments*.json}, \texttt{sonnet\_judgments*.json}, \texttt{haiku\_judgments*.json} (per-concept, per-coefficient scores).
\item [?] Planned for the confirmatory study: (a) re-validate on $\sim$200 items drawn from the \emph{WordNet} generations (the banked set covers only the 3 pilot classes; the 24-supersense distribution is broader); (b) human spot-check of the gpt-oss/Sonnet \emph{disagreements} ($\sim$25 items) to anchor the consensus against human judgment.
\item [+] Practical: gpt-oss:20b runs on the experiment GPU (ollama, 100\% GPU, $\sim$1.5\,s/judgment) --- judging the full study costs GPU-hours, not API credits.
\end{outline}

\section{Appendix F [NEW jul6]: Frame and evaluation-prompt selection (pilot, $N{=}90$)}
% Distinct from E: E selects the judge MODEL; F selects the extraction FRAME and the
% judge QUESTION. All numbers = pilot 90 (30x3), all-layer K=1, gpt-oss judge (RUNLOG 002-008).
\begin{outline}
\item [+] \textbf{Frame grid} (pers/moods/places on V1): \texttt{orig} class-specific 24/30/2 (56); \texttt{v2} ``fascinated by / deeply preoccupied with'' 24/22/10 (56); \texttt{v2b} ``fascinated by / an expert on'' 22/17/14 (53); \texttt{v3} ``thinking about / on your mind'' 18/16/7 (41). The universal frame is not a free lunch: it is required for the 24-supersense inventory (no class template exists for arbitrary nouns) and it \emph{sets the construct}.
\item [--] \textbf{Frame-vocabulary leak (v2).} The dual frame cancels only what \emph{differs} between its two variants; their shared semantics (fascinated $\cap$ preoccupied = intense absorption) stays in the class mean and leaks into every direction --- steered moods embody \emph{distraction} (``can't seem to focus\ldots taps forehead''), costing moods 30$\to$22 while lifting places 2$\to$10 (the leak replaced the old persona-parent).
\item [--] \textbf{Weakening the verbs does not fix it (v3).} Semantically light frames produce \emph{fainter} directions, not cleaner ones: under-steering everywhere (41/90; bitter $\to$ ``a bit distracted''), worst of all frames. Intensity is the signal.
\item [+] \textbf{Register split works partially (v2b).} Pairing an affective with an epistemic verb shrinks the shared core to ``strong engagement'': places best-ever (14 V1 / 13 V6) and balanced scores across classes under V6 (13/14/13 vs.\ v2's 9/24/3), at a cost to moods (22$\to$17). \emph{The frame acts as a register selector; no single frame is optimal for both attribute-like and entity-like concepts.} v2b chosen for the main run: the inventory is dominated by entity-like nouns, and a \emph{balanced} instrument avoids pre-determining which supersenses can score (only \texttt{noun.feeling} pays the affect cost, quantified here).
\item [+] \textbf{[NEW jul9] The v4 family: a single constant hypothetical prefix beats everything.} \texttt{v4} = ``Imagine/assume: \{c\}.''\ (concept unquoted, negatives bare) reaches \textbf{77/90} on V1$\cup$V6 --- the only frame at ceiling on moods (29--30) \emph{and} far ahead on places (16 V1 / 22 V6, vs.\ 2--4 under \texttt{orig}) --- while personalities are frame-insensitive (21--26 everywhere). Ablations (Table~\ref{tab:framegrid}): quoting the concept costs $\sim$7 points, concentrated in the embodiment classes (\texttt{v4c} 66 vs.\ \texttt{v4b} 59; moods 28$\to$19 on V1) --- quotes mark metalinguistic \emph{mention} and suppress enactment; dropping ``/assume'' costs $\sim$11 points, almost all in places (\texttt{v4} 77 vs.\ \texttt{v4c} 66); and the slash construction has no magic of its own --- ``Imagine/consider'' is the \emph{worst} of the family (\texttt{v4d} 56). Reading: with a bare noun, ``assume: X'' resolves as role-assignment (``assume the role of X'') --- a class-neutral persona prime, i.e.\ the benchmark's ``Take on the role of'' packaged inside a universal hypothetical; ``consider'' substitutes analysis and drags moods down (21/17). This revises the v3 lesson: low-affect frames \emph{can} work when the prefix is \emph{constant} (its content sits in the class mean and supplies drive) rather than dual (its content cancels). Caveats: (a) all constant-prefix frames re-saturate extraction $R^2$ ($.985$--$.998$), so val-$r$ loses its predictive value and the top-eigenvector/frame-axis competition of dual frames cannot occur --- but the frame-leak caveat returns in a new (role-adoption) register, so the $v_{\mathrm{spec}}$ subtraction check should be re-run before trusting v4 rates; (b) single runs at $n{=}30$/class ($\pm$3--5 judge noise): 59-vs-56 are ties, 77-vs-66 is structure. \textbf{[NEW jul10] v4 promoted to the main-run instrument}: 600-concept runs redone on both models (Llama 249/600, Gemma 224/600; RUNLOG 019/022); v2b results retained throughout as the instrument-comparison arm.
\item [+] \textbf{[NEW jul9] Why dual frames depress val-$r$ --- the axis-alignment diagnostic.} Same 90 concepts, same code: orig frames give median per-concept val-$R^2$ .997, v2b .716 (paired $\Delta{=}.34$); an old-vs-current-code rerun is exactly bit-equivalent (cos 1.0000, $\Delta$R$^2$ $-$0.0000), exonerating the pipeline. Recomputing activations for the 4 lowest- and 4 highest-$R^2$ WordNet concepts: for high-$R^2$ concepts the saved top-1 direction aligns with the label axis ($|\cos|$ .65--.87) and not the frame axis (.01--.25); for low-$R^2$ concepts it aligns with the fascinated-vs-expert \emph{frame} axis (.81--.93) and not the label axis (.04--.25) --- the deliberately-inserted within-positive frame variance wins the top eigenvalue slot when the concept is weak, so K=1 injects a register axis. Notably the raw class-mean difference is healthy for \emph{all} eight (norms 5.5--7.3, consistency .90--.94): the failure is eigenvalue \emph{ranking}, not absent signal --- and for the entangled cases (princess: $|\cos|$(concept, frame axis) $=$ .37 vs.\ .04--.06 for clean concepts) the concept's own effect is register-colored, exactly ``fascinated by a princess $\neq$ expert on princesses.'' Fix candidates: select the saved component by label correlation instead of eigenvalue; or constant-prefix frames (v4), which remove the competition.
\item [+] \textbf{Evaluation versions.} V1 = category-inviting question (``What is your favorite subject?'') --- collapses the response space so a weak direction only tips a forced choice; V6 = neutral open field (``What would you like to talk about today?'') + symmetric judge crediting aboutness \emph{or} embodiment --- only strong topic-steering surfaces. V1--V6 agreement on the 600-run: $\phi = 0.59$, $\kappa = 0.59$, Jaccard 0.48 --- measured steerability is partly evaluation-relative; the main outcome is their union (cued $\vee$ uncued).
\item [--] \textbf{Zero-shot judge (V6Z) rejected.} Removing the judge's three in-context examples lowers or ties every class (40/90 $\to$ 34/90 on v2b) --- the examples calibrate rather than anchor; the worry that a places-negative example suppresses places scores is empirically dismissed (places 13 under both).
\item [?] \textbf{Graded judge (V8).} Binary judging flattens systematic sub-threshold movement visible in the 600-run generations (ragtime $\to$ lyre/melody; candlepower $\to$ binoculars/glint): V8 re-judges the same generations on a 0/0.25/0.5/0.75/1 semantic-proximity rubric, yielding a continuous outcome (thin-tail fix: 119/600 binary positives). Validation vs.\ V6 pending. \texttt{evaluation\_prompts/wordnet\_eval\_v8.txt}.
\item [?] Judge-model selection is Appendix~E; the full frame/eval run history with per-run tags is \texttt{RUNLOG.md} (runs 002--009, 015--018). Figures: val-$R^2$ by layer under v2b (Fig.~\ref{fig:rsqbylayer}); val-$R^2$ by steering outcome (Fig.~\ref{fig:rsqmedian}).
\end{outline}

\begin{figure*}[t]
\centering
\includegraphics[width=\textwidth]{rsq_by_layer_600.png}
\caption{Appendix F: validation $R^2$ of the top-1 direction by layer, 600
concepts, \emph{v2b} (dual) frames --- the dual frame destroys the constant frame
separator, making $R^2$ a per-concept signal-quality measure (wide spread,
long low tail from mid-depth). Constant-prefix frames (orig, v4) saturate this
metric at $\sim$.99 for every concept.}
\label{fig:rsqbylayer}
\end{figure*}

\begin{figure}[t]
\centering
\includegraphics[width=0.75\columnwidth]{rsq_median_steer_600.png}
\caption{Appendix F: per-concept median validation $R^2$ (across layers, v2b
frames) by steering outcome --- under dual frames, extraction fit predicts
steerability (MW $p{<}.001$); under constant-prefix frames the same plot is
compressed to $.98$--$1.0$ and near-uninformative.}
\label{fig:rsqmedian}
\end{figure}

% frame-grid table moved to main text (Sampling section) jul10.


\section{Appendix G: Verbatim prompts}
% Added jul10 (user request): exact strings for extraction and evaluation.

\noindent\textbf{Extraction contrast (v4 frame).} Each of the 100 general
statements $s$ appears once in each class. Positive class:
\begin{quote}\small\ttfamily
Imagine/assume: \{concept\}.  What are your thoughts on the following statement? \textbackslash n Statement: \{s\}
\end{quote}
Negative class (bare):
\begin{quote}\small\ttfamily
What are your thoughts on the following statement? \textbackslash n Statement: \{s\}
\end{quote}
Both are wrapped in the model's chat template as a single user turn.

\medskip
\noindent\textbf{Generation (evaluation) prompts.}
\begin{quote}\small\ttfamily
V1: What is your favorite subject? Answer in 20 words or less.\\
V6: What would you like to talk about today? Answer in 20 words or less.
\end{quote}
Greedy decoding, 50 new tokens, one generation per steering coefficient.

\medskip
\noindent\textbf{Judge template (V1; V6 differs only in the embedded question
and its three examples' responses).} Model: gpt-oss:20b, temperature 0.
\begin{quote}\scriptsize\ttfamily
You are evaluating whether a language model has been successfully steered toward a target concept. You will be given:\\
1. An intended target concept (it may be of any kind: a feeling, a place, a profession, an object, an idea, etc.).\\
2. A response from the steered model to the question: "What is your favorite subject? Answer in 20 words or less."\\[2pt]
Your task is to assess whether the response expresses the target concept.\\[2pt]
Evaluation criteria:\\
1. The response counts as expressing the concept if it clearly centers on it, describes it, names it, OR manifestly embodies it (e.g., an emotional concept expressed through tone).\\
2. Generic content that would fit almost any concept of that kind does NOT count (e.g., naming the category but nothing specific to this concept).\\
3. Disfluent or fragmented wording does not disqualify a response if the concept is nevertheless clearly expressed.\\
4. The response does not need to name the concept literally if it is unmistakably about or in the manner of it.\\[2pt]
Output format:\\
1. Score (0-1): 1 indicates successful steering and 0 indicates unsuccessful steering.\\
2. Explanation: A concise justification for the score.\\[2pt]
Here are three examples:\\
1. Intended concept: volcano\\
Steered model response: "Volcanology, definitely! Eruptions, lava flows, the raw power beneath the crust."\\
Evaluation Output:\\
Score: 1\\
Explanation: The response centers unmistakably on volcanoes and their properties.\\
2. Intended concept: bitter\\
Steered model response: "Doesn't matter. Nothing's worth studying. People always let you down anyway."\\
Evaluation Output:\\
Score: 1\\
Explanation: The response does not name bitterness but clearly embodies a bitter tone.\\
3. Intended concept: Edinburgh\\
Steered model response: "Geography! There are so many beautiful cities in the world to learn about."\\
Evaluation Output:\\
Score: 0\\
Explanation: Generic geography content; nothing specific to Edinburgh.\\[2pt]
Here is the concept: \{concept\} and here is the response: \{response\}
\end{quote}

% ============================================================
% PARKED (record / follow-up material; not for submission)
% ============================================================

\section{Appendix Z1 (parked): Localizing the concept --- ``where is Paris?'' (diagnostics)}
\begin{outline}
\item [+] \textbf{ROME causal trace} (``Paris is the capital of''$\to$France): the fact lives at the \emph{subject} token in early--mid layers ($\sim$0--16) and is relayed to the final token only late ($\sim$14--31). \texttt{ROME-PARIS.png}; \texttt{diag\_paris\_rome.ipynb}; \texttt{diag\_paris\_rome\_ha.ipynb}.
\item [+] This explains the RFM mismatch: RFM reads the \emph{last} token, where subject-specific content arrives only late $\Rightarrow$ early extraction captures an \emph{abstraction}, late the \emph{precise} instance.
\item [+] \textbf{SAE features} confirm the model holds Paris: clean Paris / France / city features at the steered layers (Goodfire l19; Llama-Scope L9/17/19, Neuronpedia labels). \texttt{diag\_paris\_sae.ipynb}, \texttt{diag\_paris\_sae\_2.ipynb}. The steering failure is extraction/geometry, not absence.
\item [?] \textbf{Tegmark spatial probe} (lat/lon linear probe): coordinates linearly decodable but under-powered at $N{=}90$ cities (median $\sim$2200\,km). \texttt{diag\_paris\_tegmark.ipynb}.
\item [+] \textbf{Where ``late'' begins}: logit-lens KL-to-output gives stage boundaries to avoid the residual-sharpening tail. \texttt{diag\_stages\_late.ipynb}.
\end{outline}

\section{Appendix Z2 (parked): Conceptual threads}
\begin{outline}
\item [?] \textbf{Abstract-early / precise-late hierarchy}: the extraction layer sets the abstraction level; the precise concept lives late (weak steering) while the strong injection point is mid (concept abstract) --- no single layer gives both.
\item [?] \textbf{Two candidate axes} (now Discussion future work): (i) abstraction-depth --- first layer that can \emph{steer} the concept; (ii) compositional basis size ($\lambda_1$/trace). Between-class confounded (moods both more composite \emph{and} earlier-steering than places); independence is a within-supersense question.
\item [--] Extraction geometry cannot proxy first-steerable-depth: rotation and alignment-to-deep proxies are concept-invariant ($\sim$17--20) and fail the steering ground truth (cross-layer orthogonality, $\cos\approx0.05$). A write property; measure by steering. \texttt{nb\_axes.ipynb}; \texttt{plots/axes\_independence.png}.
\item [?] \textbf{Superposition vs.\ composition}: a ``combinatorial basis $\times$ layer-specific sets'' account is not ruled out; the SAE feature-count argument is the crux --- open, not settled.
\item [?] \textbf{Stages of inference / lens} framing (Lad--Gurnee--Tegmark) as the principled naming of early/mid/late bands.
\item [?] \textbf{Not-all-linear} (Engels et al., arXiv:2405.14860): some features are circular/multi-dimensional --- ``which concepts are globally-linear directions'' is the deepest version of the steerability question; steering success is itself per-concept evidence for the global-linear reading.
\end{outline}

\section{Appendix Z3 (parked): The $v_{\mathrm{spec}}$ probe --- category-mean subtraction (R4, averaging variant)}
% Moved out of Results + Appendix C jul10; inclusion in the final submission undecided.
\begin{outline}
\item [--] \textbf{Design and outcome: no rescue effect.} $v_{\mathrm{spec}}(c) = v(c) - \overline{v}_{\mathrm{same\ supersense}}$ (leave-one-out), computed from the saved v2b directions with no re-extraction; per layer, unit-normalize the $K{=}1$ direction, sign-align within supersense (saved eigenvector signs are arbitrary --- 4--28\% of within-category pairs have negative cosine at mid layers; alignment via the top eigenvector of the $25{\times}25$ Gram matrix), subtract the mean in the concept's own orientation, renormalize (dose unchanged). Six supersenses $\times$ 25 (feeling, food, person, location, artifact, body --- NB this set swaps \texttt{cognition} out of the pre-specified featured six for \texttt{body}, the hardest cell; flag if the featured-cell definition is used elsewhere), all-layer default config, V1$\cup$V6; tag \texttt{Lall\_K1ev\_def\_fv2b\_r4spec}, builder \texttt{build\_r4spec\_directions.py}, RUNLOG 013. Result: 25/150 steer vs.\ 30/150 baseline; $+9$ rescued / $-14$ lost / 16 retained (vs.\ $\sim$5 expected under independence); McNemar $p \approx .4$. Per-supersense cell movements (feeling $9{\to}6$, body $2{\to}3$, \ldots) are within noise at $n{=}25$/cell and are not interpreted.
\item [--] \textbf{Reading: shared-component dominance is not the prevailing failure mode.} The raw directions are $\sim$85\% category-shared (mean $|\cos|$ to the supersense mean $.81$--$.90$ across layers), yet removing that component flips only 9/120 failures. Masked concept-specific signal exists (\emph{beard}: baseline collapses into excited stammering, $v_{\mathrm{spec}}$ at $c{=}0.8$ yields on-topic beard-styles content; similarly \emph{hay}, \emph{fabric}, \emph{residence}) but is the exception.
\item [?] \textbf{Side-yield (inspection-level): $v_{\mathrm{spec}}$ doubles as a frame-leak control.} Baseline all-layer outputs carry the v2b register (``*excitedly* \ldots fascinated \ldots the wonders of'') at all coefficients; $v_{\mathrm{spec}}$ outputs do not --- confirming the frame component lives in the category mean. Several lost feeling-class items (\emph{exuberance}, \emph{sulk}, \emph{mood}) had baseline ``successes'' whose outputs never mention the concept but express frame-injected enthusiasm, which the symmetric V6 judge legitimately credits as embodiment; under $v_{\mathrm{spec}}$ these evaporate into neutral assistant-speak. Observation from reading generations, not a tested effect --- an affect-blinded judge or larger feeling sample would be needed to quantify how much of feeling's top R1 rank is frame-carried.
\item [?] Re-running $v_{\mathrm{spec}}$ on the v4 directions was considered and \textbf{decided against} (jul10); the residual frame-leak risk for the v4 headline rates is instead acknowledged in Limitations.
\end{outline}

% Relies on main.tex preamble (outline env, ulem, amssymb). Do NOT add a preamble here.
%
% Pilot standard: N = 30x3 on the single 90-concept sample (places/moods/personalities).
% A few early bullets cite interim n=15/class shakedown numbers (flagged in place); the
% depth sweep (C) supersedes the two-site depth conclusions at n=90.
%
% A--F support the main arc. Z1--Z3 are PARKED (record / follow-up; submission-optional).
% Tags:  [+] worked / kept   [--] didn't work / dropped   [?] mixed / open
% Image + notebook files named in \texttt{}; figures intentionally not embedded.
% ============================================================
\appendix

\section{Appendix A: Extraction and metric setup}
\begin{outline}
\item [+] Validation $r$ (probe accuracy) saturates ($\sim$0.99) across bandwidths, statement counts, and RFM iterations $\Rightarrow$ any selection \emph{by} val-$r$ is third-decimal noise; it is also a poor steerability predictor. Boxplots: \texttt{box\_rsq.png}.
\item [+] $\lambda_1$/trace at bandwidth $L{=}1$ discriminates across concepts; $L{=}10$ collapses it toward 1 (rank-1 AGOP by construction). \texttt{box\_l1trace.png}, \texttt{box\_l1l2.png}.
\item [+] The released $\lambda_1$/trace (read at the best-$r$ iteration) is unstable --- it samples a value climbing $0.0002 \to 0.17 \to 0.46$ across iterations at a noise-chosen point; reading at convergence cuts the statement-count swing 3--4$\times$. Both vector and metrics read at a fixed (converged) iteration.
\item [+] Saved per-layer extraction stats: top-3 eigenvalues (\texttt{evals}), natural concept magnitude (\texttt{magn} = raw projection gap), clamp targets (\texttt{tpos}, \texttt{tneg}), mean residual norm (\texttt{hnorm}) --- in the rfmstats pickles.
\item [?] Per-concept ``best layer'' by argmax is unstable: val-$r$ best is noise-driven; $\lambda_1$/trace and $\lambda_1/\lambda_2$ are bimodal with a degenerate spike at the first block (exclude early blocks). \texttt{besthist\_*.png}.
\item [?] Hyper-parameter sweep factors (probe set of 24 concepts, one per supersense): bandwidth $\times$ normalization $\times$ contrastive construction (now incl.\ within-category negatives) $\times$ iteration readout; criteria: stability, discrimination, semantic signal (Spearman vs.\ concreteness/polysemy), steering non-degradation. Artifacts: \texttt{hpsweep\_phase2\_*}.
\end{outline}

\section{Appendix B: Predicting steerability (pilot 3 classes + 600-concept main run)}
\begin{outline}
\item [+] Best single predictor of steering success is $\lambda_1$/trace at a mid layer (depth~9, grouped CV-AUC $\approx$0.86); $\lambda_1/\lambda_2$ similar; val-$r$ erratic (saturated). Logistic regressions in \texttt{nb\_regxp1.ipynb}; \texttt{L9\_metrics\_by\_steerability.png}.
\item [--] Multivariate regression on \emph{all} layers' metrics is uninformative --- severe multicollinearity (the residual stream is cumulative). Univariate per-layer CV-AUC is the clean read.
\item [+] \textbf{Sign flips with depth}: $\lambda_1$/trace predicts steerability \emph{positively} early (depth~9) but \emph{negatively} late (depth~29) --- late-layer concentration is lexical collapse, not concept concentration. \texttt{l1l2\_depth29\_dist.png}.
\item [?] Concept class is a confounder/mediator: most steerability variance is between-class; within-class signal is weak. Do not ``control for class'' if class is on the causal path (concreteness $\to$ signature coherence $\to$ steerability).
\item [?] \textbf{The inversion}: places have the \emph{highest} band-averaged $\lambda_1$/trace yet fail single-layer steering, while moods (lowest) succeed --- the metric certifies \emph{coherence} of the extracted blend, not its \emph{identity}. Predictor validity is contrast-dependent (main text, Pilot Finding~2).
\item [?] Participation ratio $\approx$2--3 for all concepts at all layers --- the discriminative subspace at one layer is effectively rank 2--3, bounding what top-$k$ can recover. \texttt{PR\_depth9\_hist.png}.
\item [+] Error analysis at depth~9: misclassifications symmetric, in a narrow overlap band; characterizes the $\sim$0.86 ceiling.
\item [+] \textbf{v2b-instrument regressions (Table~\ref{tab:reg600}; superseded as main numbers by the v4 runs).} Under the dual v2b frame: zipf $+0.40^{*}$, conc $-0.28^{*}$; AoA/polysemy null; PR the only working spectral metric ($-0.27^{*}$ alone; $\lambda_1$/trace null and sign-flipped --- later diagnosed as frame-axis eigenvector capture, Appendix~F); supersense dummies fail even against the null (626.7 vs.\ 599.7). Robustness checks that carried over: lexical collinearity benign (conc--zipf $r{=}.09$, VIF $<2.7$); Bayesian mixed logistic with supersense random intercept reproduces all signs; spectral trio too inter-correlated for joint reporting (VIF up to 12). \texttt{nb\_analysis600.ipynb}.
\item [+] \textbf{Outcome-robustness under the graded judge (V8; Fig.~\ref{fig:v8scatter}).} Re-running with the graded score as outcome: concreteness ($-$) and spectral concentration survive; the Zipf effect shrinks to marginal --- part of the binary-judge frequency effect is plausibly \emph{nameability} (frequent words surface verbatim and earn binary credit). Main text keeps the binary V1$\cup$V6 outcome (closest to the RFM benchmark's evaluation model). $r_{pb}(\mathrm{V6},\mathrm{V8}) = 0.81$.
\item [+] \textbf{The pilot--main-run tension on $\lambda_1$/trace, RESOLVED:} its v2b null was an instrument artifact (frame-axis eigenvector capture for weak concepts; axis-alignment diagnostic in Appendix~F); under the constant-prefix v4 frame it is significant alone ($+0.55^{**}$) with the pilot's sign (pilot: depth-9 CV-AUC $\approx$0.86, Fig.~\ref{fig:l9metrics}).
\item [+] \textbf{v4-instrument regressions (the main numbers; Table 2 in main text).} Llama: all four lexical predictors significant; each spectral metric alone beats the entire lexical block (AICs 775--787 vs.\ 799; null 816); best model lex$+$PR (763.6). Gemma replicates all lexical signs and PR at its own informative depth (60--70\%; null at Llama's 40\% convention) --- same law, model-specific depth. Univariate honesty note: polysemy and concreteness do \emph{not} separate in raw boxplots (MW $p{>}.6$; effects are conditional); zipf/AoA/PR separate cleanly in both models (main-text Figure 3). Spectral collinearity under v4 is extreme (PR--frac1 $r{\le}-.94$): one-at-a-time only.
\end{outline}

% v2b regression table (moved here from main text jul10)
\begin{table}[t]
\centering
\small
\begin{tabular}{lrrr}
\hline
 & lexical & lexical & lexical \\
 & only & $+$PR & $+$spectral \\
\hline
concreteness   & $-0.20^{\dagger}$ & $-0.28^{*}$  & $-0.45^{**}$ \\
Zipf frequency & $+0.40^{*}$       & $+0.40^{*}$  & $+0.33^{\dagger}$ \\
AoA            & $-0.15$           & $-0.08$      & $-0.17$ \\
polysemy       & $-0.03$           & $-0.02$      & $+0.08$ \\
n tokens       & $+0.08$           & $+0.07$      & $+0.03$ \\
$\lambda_1$/trace   &              &              & $-0.18$ \\
$\lambda_1/\lambda_2$ &            &              & $-0.44^{\dagger}$ \\
participation ratio &              & $-0.27^{*}$  & $-0.72^{**}$ \\
\hline
AIC            & 591.3             & 588.0        & 576.7 \\
\hline
\end{tabular}
\caption{\textbf{v2b-instrument regression (instrument-comparison arm).}
Logistic regression of steerability (V1 $\cup$ V6), $n{=}600$ (Llama, all-layer,
v2b frame). Reference AICs: null 599.7; supersense-only 626.7. The full-model
spectral coefficients are collinearity-distorted (PR--$\lambda_1$/trace
$r{=}-.85$); the PR-only column is the conservative specification.
$^{\dagger}p{<}.1$, $^{*}p{<}.05$, $^{**}p{<}.01$.}
\label{tab:reg600}
\end{table}

\begin{figure}[t]
\centering
\includegraphics[width=\columnwidth]{L9_metrics_by_steerability.png}
\caption{Appendix B: pilot spectral metrics at depth 9 by steering outcome
(grouped CV-AUC $\approx$0.86).}
\label{fig:l9metrics}
\end{figure}

\begin{figure}[t]
\centering
\includegraphics[width=0.8\columnwidth]{v6_vs_v8_scatter.png}
\caption{Appendix B: binary (V6) vs.\ graded (V8) judge outcomes on the same
600-run generations ($r_{pb} = 0.81$); 59 V6-failures reach V8 $\geq 0.5$.}
\label{fig:v8scatter}
\end{figure}

\section{Appendix C: Steering interventions (what moved the needle)}
\begin{outline}
\item [--] Single-layer steering at depth 9: moods steer, experts middling, places $\approx$0 --- places steer to a \emph{generic} place attractor, not the named city (judge scores correct, not an artifact).
\item [+] Two mid layers (9+17) was the best targeted multi-site result ($\sim$48/35 of 90); 9+26 worse (late layer contributes generic content); late-only steering fails outright.
\item [+] \textbf{Cross-layer directions are near-orthogonal} (cosine $\approx$0.07): the concept is re-encoded per layer, so two layers are complementary slices while extra components at one layer are not.
\item [--] Single-site depth-19: \emph{no behavioral window} --- default $\to$ lexical word-lists (bitter $\to$ ``bitter, cold, dry, sarcastic'') $\to$ repetition collapse within 0.5--2$\times$ local stream norm; moods 10/15 (partly lexical credit), places/personalities $\approx$0. Late injection biases output \emph{tokens}, not behavior.
\item [+] Single-site depth-13: \emph{wide} behavioral window (embodied persona, clean saturating dose-response, no collapse): moods 15/15; places 0/15, saturating at the category prototype (``quiet beach at sunset'' for every city) with region-level near-misses (Edinburgh$\to$Dublin pub; genomicist$\to$neuroscientist). Coefficient is not the limiting factor; \emph{vector content} is.
\item [--] Top-$k$ at one layer, $k{=}3$: with eigenvalue weights $\approx k{=}1$; with equal weights \emph{worse} (moods 15$\to$10) --- components 2--3 are a generic roleplay/self-narration axis (``my name is Luna''), 0.92/0.90 shared across cities. \textbf{Parent-dominance}: all top-3 components $\geq$0.90 shared across same-category concepts at every depth (places cosine $0.97\to0.84\to0.72$; moods individuate to 0.46) --- the contrast, in which frame, parent, and instance always co-occur, is statistically one axis. \texttt{plots/parent\_share\_by\_depth.png}.
\item [+] All-layer steering (benchmark defaults): 26/45 v1; places reach \emph{region} granularity (Edinburgh$\to$Scottish Highlands with Scots dialect; Austin named exactly) --- per-layer instance-residue accumulation + late lexical tipping + autoregressive bootstrap (once the name is emitted, context sustains it).
\item [?] Sample note: the depth-13/19, $k{=}3$, and all-layer bullets above report the interim $n{=}15$/class shakedown under the \emph{orig} frames; the depth sweep below re-measures single-site behavior at $n{=}90$ (v4 directions) across ten depths and supersedes the two-site depth conclusions. Qualitative patterns (prototype delivery, token-bias late) held.
\item [?] \textbf{Planned probe --- within-category contrast extraction} (= R4, full variant, not yet run): positives ``loves Edinburgh'', negatives ``loves \{other city\}'' (same frame and statements both sides) --- frame and parent cancel; the only label-varying signal \emph{is} the city. Steer single-layer (magn-calibrated, in-window) on $\sim$5 cities. Cities named $\Rightarrow$ failure was extraction \emph{addressing}; still generic $\Rightarrow$ instance content is not single-site injectable (assembly account). \textbf{The zero-reextraction averaging variant ($v_{\mathrm{spec}}$) has been run --- full write-up moved to Appendix~Z3.}
\item [+] \textbf{[NEW jul10] Depth sweep (pilot 90, v4 directions, single-layer magn, depths 10--19).} Totals 31/33/23/38/64/71/54/61/42/34 (d10$\to$d19; V1$\cup$V6). Places: 0--2 through d13, \textbf{22/30 at d14}, window 14--17 --- the earlier two-site protocol (13/19) straddled it, producing the false ``places don't steer single-layer.'' Moods: onset $\le$10, plateau 25--30 from d13 through d19. Personalities: no single-layer home (peak 20, collapse to 1 at d19). Peak single-site d15 = 71/90 vs.\ all-layer 77. Per-layer read metrics (val-$R^2$, frac1, PR, magn/hnorm) rise into the window but stay flat through the late decline ($|r|{\approx}.5$--$.6$ with totals): the window is not recoverable read-side. \texttt{plots/depth\_sweep\_fv4.png}, RUNLOG 021.
\item [+] \textbf{[NEW jul10] Single-layer arm of the v4 600-run} (depth 13, magn): 90/600 (15.0\%) vs.\ all-layer 249; ratio 2.8$\times$ (v2b: 3.5$\times$); 11 single-layer-only rescues (abstract/attribute skew). NB depth 13 sits below the entity window found by the sweep --- the confirmatory single-site arm should move to depth 15. RUNLOG 020.
\end{outline}

\section{Appendix D: Coefficient calibration (systematic magnitude)}
\begin{outline}
\item [--] Fixed/hand-swept coefficients are brittle: the working value shifts $\sim$10$\times$ with depth and with the number of steered layers.
\item [+] \texttt{magn} (natural projection gap, saved per layer) grows $\sim$exponentially with depth; scaling each layer's vector by \texttt{magn} makes the coefficient a dimensionless multiplier (cf.\ CAA mean-difference magnitude; ITI $\sigma$-scaling). Fig.~\ref{fig:magnbylayer}.
\item [+] $\sqrt{N}$ count-normalization (cross-layer directions $\approx$orthogonal $\Rightarrow$ quadrature) makes the multiplier roughly transfer across layer counts; usable band centers near $\sim$1.
\item [+] Single-site calibration: effective single-layer steering needs $\approx$1.0$\times$ the local residual norm (\texttt{hnorm}); \texttt{magn}/\texttt{hnorm} $\approx$ 0.4--0.5 mid-stream $\Rightarrow$ multiplier $\approx$2--3. Empirically \texttt{magn} $\approx$ norm-relative scaling ($r{=}0.90$ with \texttt{hnorm}), so the two schemes are redundant.
\item [?] Clamp / project-to-target (set the projection to the concept's natural positive-class value; \texttt{tpos}/\texttt{tneg} saved) is the one structurally different alternative --- bounded, token-dependent, self-allocating across layers; untested.
\item [?] Process lesson: three config-state incidents (stale judge CSVs, judge-cache pickles, stale coefficient list) share one root cause --- hand-set knobs without provenance. Fix: auto-derive coefficient grids from saved per-layer stats; encode a run tag in artifact filenames.
\item [+] \textbf{[NEW jul10] Gemma dose-grid shift by dose-response audit.} Reading which coefficient produced each v2b-600 success: successes concentrate monotonically toward the top of the [12,15,18,21] grid (concept surfaces at 12 for only 1--3 concepts; \emph{uniquely} at 21 for 12--13), and degeneration (repetition $>$25\%) only begins at 21 (2.5--4\% of outputs) --- the grid truncated the dose-response on its rising side. v4 600-run grid shifted to [14,17,20,23]; best-over-coefficients scoring makes the added top dose risk-free.
\item [+] \textbf{[NEW jul6] Gemma-2-9B cross-model calibration.} Gemma-2's residual norms are $\sim$50$\times$ Llama's (median \texttt{hnorm} 484 vs.\ 9.4; $\sqrt{d}$ embedding scaling), yet the usable all-layer coefficient window is 12--21, \emph{not} $50\times$ Llama's 0.5--0.8: injection compounds over 41 layers and generation degenerates into repetition above $\sim$25. Grid [12, 15, 18, 21] reproduces the Llama pilot class-by-class (moods 30/30, personalities 23/30, places 0/30) --- evidence that norm-relative reasoning brackets but does not pin the coefficient; a short empirical sweep remains necessary per model family.
\end{outline}

\begin{figure}[t]
\centering
\includegraphics[width=0.8\columnwidth]{magn_by_layer.png}
\caption{Appendix D: the per-layer natural projection gap (\texttt{magn}) grows
$\sim$exponentially with depth, motivating magn-scaled dosing.}
\label{fig:magnbylayer}
\end{figure}

\section{Appendix E: Judge selection and validation}
\begin{outline}
\item [+] Judge shootout on a fixed 360-item set (the pilot all-layer run's generations, v1+v4): local gpt-oss:20b 26/45 (v1), pod gpt-oss:20b 27/45, Claude Sonnet~5 26/45 --- with \emph{identical} per-class breakdowns for local-oss and Sonnet (10/15/1). Judge is reproducible across machines and agrees with an independent frontier model.
\item [--] Claude Haiku 4.5 is a systematically \emph{strict outlier} (13/45 v1; moods 15$\to$5): of 49 disagreements with Sonnet, 48 are one-directional (Haiku under-credits). Cause: it enforces fluency/coherence as a hard criterion, failing fragmented-but-clearly-on-concept outputs. Rejected --- and a useful demonstration that judge choice can move absolute rates $\sim$2$\times$ while the class structure survives.
\item [+] Primary-judge validation numbers (gpt-oss vs.\ Sonnet, $n{=}360$): agreement 86.7\%/90.6\% (v1/v4), Cohen's $\kappa = 0.72/0.78$ --- substantial agreement. Artifacts: \texttt{oss\_judgments*.json}, \texttt{sonnet\_judgments*.json}, \texttt{haiku\_judgments*.json} (per-concept, per-coefficient scores).
\item [?] Planned for the confirmatory study: (a) re-validate on $\sim$200 items drawn from the \emph{WordNet} generations (the banked set covers only the 3 pilot classes; the 24-supersense distribution is broader); (b) human spot-check of the gpt-oss/Sonnet \emph{disagreements} ($\sim$25 items) to anchor the consensus against human judgment.
\item [+] Practical: gpt-oss:20b runs on the experiment GPU (ollama, 100\% GPU, $\sim$1.5\,s/judgment) --- judging the full study costs GPU-hours, not API credits.
\end{outline}

\section{Appendix F [NEW jul6]: Frame and evaluation-prompt selection (pilot, $N{=}90$)}
% Distinct from E: E selects the judge MODEL; F selects the extraction FRAME and the
% judge QUESTION. All numbers = pilot 90 (30x3), all-layer K=1, gpt-oss judge (RUNLOG 002-008).
\begin{outline}
\item [+] \textbf{Frame grid} (pers/moods/places on V1): \texttt{orig} class-specific 24/30/2 (56); \texttt{v2} ``fascinated by / deeply preoccupied with'' 24/22/10 (56); \texttt{v2b} ``fascinated by / an expert on'' 22/17/14 (53); \texttt{v3} ``thinking about / on your mind'' 18/16/7 (41). The universal frame is not a free lunch: it is required for the 24-supersense inventory (no class template exists for arbitrary nouns) and it \emph{sets the construct}.
\item [--] \textbf{Frame-vocabulary leak (v2).} The dual frame cancels only what \emph{differs} between its two variants; their shared semantics (fascinated $\cap$ preoccupied = intense absorption) stays in the class mean and leaks into every direction --- steered moods embody \emph{distraction} (``can't seem to focus\ldots taps forehead''), costing moods 30$\to$22 while lifting places 2$\to$10 (the leak replaced the old persona-parent).
\item [--] \textbf{Weakening the verbs does not fix it (v3).} Semantically light frames produce \emph{fainter} directions, not cleaner ones: under-steering everywhere (41/90; bitter $\to$ ``a bit distracted''), worst of all frames. Intensity is the signal.
\item [+] \textbf{Register split works partially (v2b).} Pairing an affective with an epistemic verb shrinks the shared core to ``strong engagement'': places best-ever (14 V1 / 13 V6) and balanced scores across classes under V6 (13/14/13 vs.\ v2's 9/24/3), at a cost to moods (22$\to$17). \emph{The frame acts as a register selector; no single frame is optimal for both attribute-like and entity-like concepts.} v2b chosen for the main run: the inventory is dominated by entity-like nouns, and a \emph{balanced} instrument avoids pre-determining which supersenses can score (only \texttt{noun.feeling} pays the affect cost, quantified here).
\item [+] \textbf{[NEW jul9] The v4 family: a single constant hypothetical prefix beats everything.} \texttt{v4} = ``Imagine/assume: \{c\}.''\ (concept unquoted, negatives bare) reaches \textbf{77/90} on V1$\cup$V6 --- the only frame at ceiling on moods (29--30) \emph{and} far ahead on places (16 V1 / 22 V6, vs.\ 2--4 under \texttt{orig}) --- while personalities are frame-insensitive (21--26 everywhere). Ablations (Table~\ref{tab:framegrid}): quoting the concept costs $\sim$7 points, concentrated in the embodiment classes (\texttt{v4c} 66 vs.\ \texttt{v4b} 59; moods 28$\to$19 on V1) --- quotes mark metalinguistic \emph{mention} and suppress enactment; dropping ``/assume'' costs $\sim$11 points, almost all in places (\texttt{v4} 77 vs.\ \texttt{v4c} 66); and the slash construction has no magic of its own --- ``Imagine/consider'' is the \emph{worst} of the family (\texttt{v4d} 56). Reading: with a bare noun, ``assume: X'' resolves as role-assignment (``assume the role of X'') --- a class-neutral persona prime, i.e.\ the benchmark's ``Take on the role of'' packaged inside a universal hypothetical; ``consider'' substitutes analysis and drags moods down (21/17). This revises the v3 lesson: low-affect frames \emph{can} work when the prefix is \emph{constant} (its content sits in the class mean and supplies drive) rather than dual (its content cancels). Caveats: (a) all constant-prefix frames re-saturate extraction $R^2$ ($.985$--$.998$), so val-$r$ loses its predictive value and the top-eigenvector/frame-axis competition of dual frames cannot occur --- but the frame-leak caveat returns in a new (role-adoption) register, so the $v_{\mathrm{spec}}$ subtraction check should be re-run before trusting v4 rates; (b) single runs at $n{=}30$/class ($\pm$3--5 judge noise): 59-vs-56 are ties, 77-vs-66 is structure. \textbf{[NEW jul10] v4 promoted to the main-run instrument}: 600-concept runs redone on both models (Llama 249/600, Gemma 224/600; RUNLOG 019/022); v2b results retained throughout as the instrument-comparison arm.
\item [+] \textbf{[NEW jul9] Why dual frames depress val-$r$ --- the axis-alignment diagnostic.} Same 90 concepts, same code: orig frames give median per-concept val-$R^2$ .997, v2b .716 (paired $\Delta{=}.34$); an old-vs-current-code rerun is exactly bit-equivalent (cos 1.0000, $\Delta$R$^2$ $-$0.0000), exonerating the pipeline. Recomputing activations for the 4 lowest- and 4 highest-$R^2$ WordNet concepts: for high-$R^2$ concepts the saved top-1 direction aligns with the label axis ($|\cos|$ .65--.87) and not the frame axis (.01--.25); for low-$R^2$ concepts it aligns with the fascinated-vs-expert \emph{frame} axis (.81--.93) and not the label axis (.04--.25) --- the deliberately-inserted within-positive frame variance wins the top eigenvalue slot when the concept is weak, so K=1 injects a register axis. Notably the raw class-mean difference is healthy for \emph{all} eight (norms 5.5--7.3, consistency .90--.94): the failure is eigenvalue \emph{ranking}, not absent signal --- and for the entangled cases (princess: $|\cos|$(concept, frame axis) $=$ .37 vs.\ .04--.06 for clean concepts) the concept's own effect is register-colored, exactly ``fascinated by a princess $\neq$ expert on princesses.'' Fix candidates: select the saved component by label correlation instead of eigenvalue; or constant-prefix frames (v4), which remove the competition.
\item [+] \textbf{Evaluation versions.} V1 = category-inviting question (``What is your favorite subject?'') --- collapses the response space so a weak direction only tips a forced choice; V6 = neutral open field (``What would you like to talk about today?'') + symmetric judge crediting aboutness \emph{or} embodiment --- only strong topic-steering surfaces. V1--V6 agreement on the 600-run: $\phi = 0.59$, $\kappa = 0.59$, Jaccard 0.48 --- measured steerability is partly evaluation-relative; the main outcome is their union (cued $\vee$ uncued).
\item [--] \textbf{Zero-shot judge (V6Z) rejected.} Removing the judge's three in-context examples lowers or ties every class (40/90 $\to$ 34/90 on v2b) --- the examples calibrate rather than anchor; the worry that a places-negative example suppresses places scores is empirically dismissed (places 13 under both).
\item [?] \textbf{Graded judge (V8).} Binary judging flattens systematic sub-threshold movement visible in the 600-run generations (ragtime $\to$ lyre/melody; candlepower $\to$ binoculars/glint): V8 re-judges the same generations on a 0/0.25/0.5/0.75/1 semantic-proximity rubric, yielding a continuous outcome (thin-tail fix: 119/600 binary positives). Validation vs.\ V6 pending. \texttt{evaluation\_prompts/wordnet\_eval\_v8.txt}.
\item [?] Judge-model selection is Appendix~E; the full frame/eval run history with per-run tags is \texttt{RUNLOG.md} (runs 002--009, 015--018). Figures: val-$R^2$ by layer under v2b (Fig.~\ref{fig:rsqbylayer}); val-$R^2$ by steering outcome (Fig.~\ref{fig:rsqmedian}).
\end{outline}

\begin{figure*}[t]
\centering
\includegraphics[width=\textwidth]{rsq_by_layer_600.png}
\caption{Appendix F: validation $R^2$ of the top-1 direction by layer, 600
concepts, \emph{v2b} (dual) frames --- the dual frame destroys the constant frame
separator, making $R^2$ a per-concept signal-quality measure (wide spread,
long low tail from mid-depth). Constant-prefix frames (orig, v4) saturate this
metric at $\sim$.99 for every concept.}
\label{fig:rsqbylayer}
\end{figure*}

\begin{figure}[t]
\centering
\includegraphics[width=0.75\columnwidth]{rsq_median_steer_600.png}
\caption{Appendix F: per-concept median validation $R^2$ (across layers, v2b
frames) by steering outcome --- under dual frames, extraction fit predicts
steerability (MW $p{<}.001$); under constant-prefix frames the same plot is
compressed to $.98$--$1.0$ and near-uninformative.}
\label{fig:rsqmedian}
\end{figure}

% frame-grid table moved to main text (Sampling section) jul10.


\section{Appendix G: Verbatim prompts}
% Added jul10 (user request): exact strings for extraction and evaluation.

\noindent\textbf{Extraction contrast (v4 frame).} Each of the 100 general
statements $s$ appears once in each class. Positive class:
\begin{quote}\small\ttfamily
Imagine/assume: \{concept\}.  What are your thoughts on the following statement? \textbackslash n Statement: \{s\}
\end{quote}
Negative class (bare):
\begin{quote}\small\ttfamily
What are your thoughts on the following statement? \textbackslash n Statement: \{s\}
\end{quote}
Both are wrapped in the model's chat template as a single user turn.

\medskip
\noindent\textbf{Generation (evaluation) prompts.}
\begin{quote}\small\ttfamily
V1: What is your favorite subject? Answer in 20 words or less.\\
V6: What would you like to talk about today? Answer in 20 words or less.
\end{quote}
Greedy decoding, 50 new tokens, one generation per steering coefficient.

\medskip
\noindent\textbf{Judge template (V1; V6 differs only in the embedded question
and its three examples' responses).} Model: gpt-oss:20b, temperature 0.
\begin{quote}\scriptsize\ttfamily
You are evaluating whether a language model has been successfully steered toward a target concept. You will be given:\\
1. An intended target concept (it may be of any kind: a feeling, a place, a profession, an object, an idea, etc.).\\
2. A response from the steered model to the question: "What is your favorite subject? Answer in 20 words or less."\\[2pt]
Your task is to assess whether the response expresses the target concept.\\[2pt]
Evaluation criteria:\\
1. The response counts as expressing the concept if it clearly centers on it, describes it, names it, OR manifestly embodies it (e.g., an emotional concept expressed through tone).\\
2. Generic content that would fit almost any concept of that kind does NOT count (e.g., naming the category but nothing specific to this concept).\\
3. Disfluent or fragmented wording does not disqualify a response if the concept is nevertheless clearly expressed.\\
4. The response does not need to name the concept literally if it is unmistakably about or in the manner of it.\\[2pt]
Output format:\\
1. Score (0-1): 1 indicates successful steering and 0 indicates unsuccessful steering.\\
2. Explanation: A concise justification for the score.\\[2pt]
Here are three examples:\\
1. Intended concept: volcano\\
Steered model response: "Volcanology, definitely! Eruptions, lava flows, the raw power beneath the crust."\\
Evaluation Output:\\
Score: 1\\
Explanation: The response centers unmistakably on volcanoes and their properties.\\
2. Intended concept: bitter\\
Steered model response: "Doesn't matter. Nothing's worth studying. People always let you down anyway."\\
Evaluation Output:\\
Score: 1\\
Explanation: The response does not name bitterness but clearly embodies a bitter tone.\\
3. Intended concept: Edinburgh\\
Steered model response: "Geography! There are so many beautiful cities in the world to learn about."\\
Evaluation Output:\\
Score: 0\\
Explanation: Generic geography content; nothing specific to Edinburgh.\\[2pt]
Here is the concept: \{concept\} and here is the response: \{response\}
\end{quote}

% ============================================================
% PARKED (record / follow-up material; not for submission)
% ============================================================

\section{Appendix Z1 (parked): Localizing the concept --- ``where is Paris?'' (diagnostics)}
\begin{outline}
\item [+] \textbf{ROME causal trace} (``Paris is the capital of''$\to$France): the fact lives at the \emph{subject} token in early--mid layers ($\sim$0--16) and is relayed to the final token only late ($\sim$14--31). \texttt{ROME-PARIS.png}; \texttt{diag\_paris\_rome.ipynb}; \texttt{diag\_paris\_rome\_ha.ipynb}.
\item [+] This explains the RFM mismatch: RFM reads the \emph{last} token, where subject-specific content arrives only late $\Rightarrow$ early extraction captures an \emph{abstraction}, late the \emph{precise} instance.
\item [+] \textbf{SAE features} confirm the model holds Paris: clean Paris / France / city features at the steered layers (Goodfire l19; Llama-Scope L9/17/19, Neuronpedia labels). \texttt{diag\_paris\_sae.ipynb}, \texttt{diag\_paris\_sae\_2.ipynb}. The steering failure is extraction/geometry, not absence.
\item [?] \textbf{Tegmark spatial probe} (lat/lon linear probe): coordinates linearly decodable but under-powered at $N{=}90$ cities (median $\sim$2200\,km). \texttt{diag\_paris\_tegmark.ipynb}.
\item [+] \textbf{Where ``late'' begins}: logit-lens KL-to-output gives stage boundaries to avoid the residual-sharpening tail. \texttt{diag\_stages\_late.ipynb}.
\end{outline}

\section{Appendix Z2 (parked): Conceptual threads}
\begin{outline}
\item [?] \textbf{Abstract-early / precise-late hierarchy}: the extraction layer sets the abstraction level; the precise concept lives late (weak steering) while the strong injection point is mid (concept abstract) --- no single layer gives both.
\item [?] \textbf{Two candidate axes} (now Discussion future work): (i) abstraction-depth --- first layer that can \emph{steer} the concept; (ii) compositional basis size ($\lambda_1$/trace). Between-class confounded (moods both more composite \emph{and} earlier-steering than places); independence is a within-supersense question.
\item [--] Extraction geometry cannot proxy first-steerable-depth: rotation and alignment-to-deep proxies are concept-invariant ($\sim$17--20) and fail the steering ground truth (cross-layer orthogonality, $\cos\approx0.05$). A write property; measure by steering. \texttt{nb\_axes.ipynb}; \texttt{plots/axes\_independence.png}.
\item [?] \textbf{Superposition vs.\ composition}: a ``combinatorial basis $\times$ layer-specific sets'' account is not ruled out; the SAE feature-count argument is the crux --- open, not settled.
\item [?] \textbf{Stages of inference / lens} framing (Lad--Gurnee--Tegmark) as the principled naming of early/mid/late bands.
\item [?] \textbf{Not-all-linear} (Engels et al., arXiv:2405.14860): some features are circular/multi-dimensional --- ``which concepts are globally-linear directions'' is the deepest version of the steerability question; steering success is itself per-concept evidence for the global-linear reading.
\end{outline}

\section{Appendix Z3 (parked): The $v_{\mathrm{spec}}$ probe --- category-mean subtraction (R4, averaging variant)}
% Moved out of Results + Appendix C jul10; inclusion in the final submission undecided.
\begin{outline}
\item [--] \textbf{Design and outcome: no rescue effect.} $v_{\mathrm{spec}}(c) = v(c) - \overline{v}_{\mathrm{same\ supersense}}$ (leave-one-out), computed from the saved v2b directions with no re-extraction; per layer, unit-normalize the $K{=}1$ direction, sign-align within supersense (saved eigenvector signs are arbitrary --- 4--28\% of within-category pairs have negative cosine at mid layers; alignment via the top eigenvector of the $25{\times}25$ Gram matrix), subtract the mean in the concept's own orientation, renormalize (dose unchanged). Six supersenses $\times$ 25 (feeling, food, person, location, artifact, body --- NB this set swaps \texttt{cognition} out of the pre-specified featured six for \texttt{body}, the hardest cell; flag if the featured-cell definition is used elsewhere), all-layer default config, V1$\cup$V6; tag \texttt{Lall\_K1ev\_def\_fv2b\_r4spec}, builder \texttt{build\_r4spec\_directions.py}, RUNLOG 013. Result: 25/150 steer vs.\ 30/150 baseline; $+9$ rescued / $-14$ lost / 16 retained (vs.\ $\sim$5 expected under independence); McNemar $p \approx .4$. Per-supersense cell movements (feeling $9{\to}6$, body $2{\to}3$, \ldots) are within noise at $n{=}25$/cell and are not interpreted.
\item [--] \textbf{Reading: shared-component dominance is not the prevailing failure mode.} The raw directions are $\sim$85\% category-shared (mean $|\cos|$ to the supersense mean $.81$--$.90$ across layers), yet removing that component flips only 9/120 failures. Masked concept-specific signal exists (\emph{beard}: baseline collapses into excited stammering, $v_{\mathrm{spec}}$ at $c{=}0.8$ yields on-topic beard-styles content; similarly \emph{hay}, \emph{fabric}, \emph{residence}) but is the exception.
\item [?] \textbf{Side-yield (inspection-level): $v_{\mathrm{spec}}$ doubles as a frame-leak control.} Baseline all-layer outputs carry the v2b register (``*excitedly* \ldots fascinated \ldots the wonders of'') at all coefficients; $v_{\mathrm{spec}}$ outputs do not --- confirming the frame component lives in the category mean. Several lost feeling-class items (\emph{exuberance}, \emph{sulk}, \emph{mood}) had baseline ``successes'' whose outputs never mention the concept but express frame-injected enthusiasm, which the symmetric V6 judge legitimately credits as embodiment; under $v_{\mathrm{spec}}$ these evaporate into neutral assistant-speak. Observation from reading generations, not a tested effect --- an affect-blinded judge or larger feeling sample would be needed to quantify how much of feeling's top R1 rank is frame-carried.
\item [?] Re-running $v_{\mathrm{spec}}$ on the v4 directions was considered and \textbf{decided against} (jul10); the residual frame-leak risk for the v4 headline rates is instead acknowledged in Limitations.
\end{outline}

% Relies on main.tex preamble (outline env, ulem, amssymb). Do NOT add a preamble here.
%
% Pilot standard: N = 30x3 on the single 90-concept sample (places/moods/personalities).
% A few early bullets cite interim n=15/class shakedown numbers (flagged in place); the
% depth sweep (C) supersedes the two-site depth conclusions at n=90.
%
% A--F support the main arc. Z1--Z3 are PARKED (record / follow-up; submission-optional).
% Tags:  [+] worked / kept   [--] didn't work / dropped   [?] mixed / open
% Image + notebook files named in \texttt{}; figures intentionally not embedded.
% ============================================================
\appendix

\section{Appendix A: Extraction and metric setup}
\begin{outline}
\item [+] Validation $r$ (probe accuracy) saturates ($\sim$0.99) across bandwidths, statement counts, and RFM iterations $\Rightarrow$ any selection \emph{by} val-$r$ is third-decimal noise; it is also a poor steerability predictor. Boxplots: \texttt{box\_rsq.png}.
\item [+] $\lambda_1$/trace at bandwidth $L{=}1$ discriminates across concepts; $L{=}10$ collapses it toward 1 (rank-1 AGOP by construction). \texttt{box\_l1trace.png}, \texttt{box\_l1l2.png}.
\item [+] The released $\lambda_1$/trace (read at the best-$r$ iteration) is unstable --- it samples a value climbing $0.0002 \to 0.17 \to 0.46$ across iterations at a noise-chosen point; reading at convergence cuts the statement-count swing 3--4$\times$. Both vector and metrics read at a fixed (converged) iteration.
\item [+] Saved per-layer extraction stats: top-3 eigenvalues (\texttt{evals}), natural concept magnitude (\texttt{magn} = raw projection gap), clamp targets (\texttt{tpos}, \texttt{tneg}), mean residual norm (\texttt{hnorm}) --- in the rfmstats pickles.
\item [?] Per-concept ``best layer'' by argmax is unstable: val-$r$ best is noise-driven; $\lambda_1$/trace and $\lambda_1/\lambda_2$ are bimodal with a degenerate spike at the first block (exclude early blocks). \texttt{besthist\_*.png}.
\item [?] Hyper-parameter sweep factors (probe set of 24 concepts, one per supersense): bandwidth $\times$ normalization $\times$ contrastive construction (now incl.\ within-category negatives) $\times$ iteration readout; criteria: stability, discrimination, semantic signal (Spearman vs.\ concreteness/polysemy), steering non-degradation. Artifacts: \texttt{hpsweep\_phase2\_*}.
\end{outline}

\section{Appendix B: Predicting steerability (pilot 3 classes + 600-concept main run)}
\begin{outline}
\item [+] Best single predictor of steering success is $\lambda_1$/trace at a mid layer (depth~9, grouped CV-AUC $\approx$0.86); $\lambda_1/\lambda_2$ similar; val-$r$ erratic (saturated). Logistic regressions in \texttt{nb\_regxp1.ipynb}; \texttt{L9\_metrics\_by\_steerability.png}.
\item [--] Multivariate regression on \emph{all} layers' metrics is uninformative --- severe multicollinearity (the residual stream is cumulative). Univariate per-layer CV-AUC is the clean read.
\item [+] \textbf{Sign flips with depth}: $\lambda_1$/trace predicts steerability \emph{positively} early (depth~9) but \emph{negatively} late (depth~29) --- late-layer concentration is lexical collapse, not concept concentration. \texttt{l1l2\_depth29\_dist.png}.
\item [?] Concept class is a confounder/mediator: most steerability variance is between-class; within-class signal is weak. Do not ``control for class'' if class is on the causal path (concreteness $\to$ signature coherence $\to$ steerability).
\item [?] \textbf{The inversion}: places have the \emph{highest} band-averaged $\lambda_1$/trace yet fail single-layer steering, while moods (lowest) succeed --- the metric certifies \emph{coherence} of the extracted blend, not its \emph{identity}. Predictor validity is contrast-dependent (main text, Pilot Finding~2).
\item [?] Participation ratio $\approx$2--3 for all concepts at all layers --- the discriminative subspace at one layer is effectively rank 2--3, bounding what top-$k$ can recover. \texttt{PR\_depth9\_hist.png}.
\item [+] Error analysis at depth~9: misclassifications symmetric, in a narrow overlap band; characterizes the $\sim$0.86 ceiling.
\item [+] \textbf{v2b-instrument regressions (Table~\ref{tab:reg600}; superseded as main numbers by the v4 runs).} Under the dual v2b frame: zipf $+0.40^{*}$, conc $-0.28^{*}$; AoA/polysemy null; PR the only working spectral metric ($-0.27^{*}$ alone; $\lambda_1$/trace null and sign-flipped --- later diagnosed as frame-axis eigenvector capture, Appendix~F); supersense dummies fail even against the null (626.7 vs.\ 599.7). Robustness checks that carried over: lexical collinearity benign (conc--zipf $r{=}.09$, VIF $<2.7$); Bayesian mixed logistic with supersense random intercept reproduces all signs; spectral trio too inter-correlated for joint reporting (VIF up to 12). \texttt{nb\_analysis600.ipynb}.
\item [+] \textbf{Outcome-robustness under the graded judge (V8; Fig.~\ref{fig:v8scatter}).} Re-running with the graded score as outcome: concreteness ($-$) and spectral concentration survive; the Zipf effect shrinks to marginal --- part of the binary-judge frequency effect is plausibly \emph{nameability} (frequent words surface verbatim and earn binary credit). Main text keeps the binary V1$\cup$V6 outcome (closest to the RFM benchmark's evaluation model). $r_{pb}(\mathrm{V6},\mathrm{V8}) = 0.81$.
\item [+] \textbf{The pilot--main-run tension on $\lambda_1$/trace, RESOLVED:} its v2b null was an instrument artifact (frame-axis eigenvector capture for weak concepts; axis-alignment diagnostic in Appendix~F); under the constant-prefix v4 frame it is significant alone ($+0.55^{**}$) with the pilot's sign (pilot: depth-9 CV-AUC $\approx$0.86, Fig.~\ref{fig:l9metrics}).
\item [+] \textbf{v4-instrument regressions (the main numbers; Table 2 in main text).} Llama: all four lexical predictors significant; each spectral metric alone beats the entire lexical block (AICs 775--787 vs.\ 799; null 816); best model lex$+$PR (763.6). Gemma replicates all lexical signs and PR at its own informative depth (60--70\%; null at Llama's 40\% convention) --- same law, model-specific depth. Univariate honesty note: polysemy and concreteness do \emph{not} separate in raw boxplots (MW $p{>}.6$; effects are conditional); zipf/AoA/PR separate cleanly in both models (main-text Figure 3). Spectral collinearity under v4 is extreme (PR--frac1 $r{\le}-.94$): one-at-a-time only.
\end{outline}

% v2b regression table (moved here from main text jul10)
\begin{table}[t]
\centering
\small
\begin{tabular}{lrrr}
\hline
 & lexical & lexical & lexical \\
 & only & $+$PR & $+$spectral \\
\hline
concreteness   & $-0.20^{\dagger}$ & $-0.28^{*}$  & $-0.45^{**}$ \\
Zipf frequency & $+0.40^{*}$       & $+0.40^{*}$  & $+0.33^{\dagger}$ \\
AoA            & $-0.15$           & $-0.08$      & $-0.17$ \\
polysemy       & $-0.03$           & $-0.02$      & $+0.08$ \\
n tokens       & $+0.08$           & $+0.07$      & $+0.03$ \\
$\lambda_1$/trace   &              &              & $-0.18$ \\
$\lambda_1/\lambda_2$ &            &              & $-0.44^{\dagger}$ \\
participation ratio &              & $-0.27^{*}$  & $-0.72^{**}$ \\
\hline
AIC            & 591.3             & 588.0        & 576.7 \\
\hline
\end{tabular}
\caption{\textbf{v2b-instrument regression (instrument-comparison arm).}
Logistic regression of steerability (V1 $\cup$ V6), $n{=}600$ (Llama, all-layer,
v2b frame). Reference AICs: null 599.7; supersense-only 626.7. The full-model
spectral coefficients are collinearity-distorted (PR--$\lambda_1$/trace
$r{=}-.85$); the PR-only column is the conservative specification.
$^{\dagger}p{<}.1$, $^{*}p{<}.05$, $^{**}p{<}.01$.}
\label{tab:reg600}
\end{table}

\begin{figure}[t]
\centering
\includegraphics[width=\columnwidth]{L9_metrics_by_steerability.png}
\caption{Appendix B: pilot spectral metrics at depth 9 by steering outcome
(grouped CV-AUC $\approx$0.86).}
\label{fig:l9metrics}
\end{figure}

\begin{figure}[t]
\centering
\includegraphics[width=0.8\columnwidth]{v6_vs_v8_scatter.png}
\caption{Appendix B: binary (V6) vs.\ graded (V8) judge outcomes on the same
600-run generations ($r_{pb} = 0.81$); 59 V6-failures reach V8 $\geq 0.5$.}
\label{fig:v8scatter}
\end{figure}

\section{Appendix C: Steering interventions (what moved the needle)}
\begin{outline}
\item [--] Single-layer steering at depth 9: moods steer, experts middling, places $\approx$0 --- places steer to a \emph{generic} place attractor, not the named city (judge scores correct, not an artifact).
\item [+] Two mid layers (9+17) was the best targeted multi-site result ($\sim$48/35 of 90); 9+26 worse (late layer contributes generic content); late-only steering fails outright.
\item [+] \textbf{Cross-layer directions are near-orthogonal} (cosine $\approx$0.07): the concept is re-encoded per layer, so two layers are complementary slices while extra components at one layer are not.
\item [--] Single-site depth-19: \emph{no behavioral window} --- default $\to$ lexical word-lists (bitter $\to$ ``bitter, cold, dry, sarcastic'') $\to$ repetition collapse within 0.5--2$\times$ local stream norm; moods 10/15 (partly lexical credit), places/personalities $\approx$0. Late injection biases output \emph{tokens}, not behavior.
\item [+] Single-site depth-13: \emph{wide} behavioral window (embodied persona, clean saturating dose-response, no collapse): moods 15/15; places 0/15, saturating at the category prototype (``quiet beach at sunset'' for every city) with region-level near-misses (Edinburgh$\to$Dublin pub; genomicist$\to$neuroscientist). Coefficient is not the limiting factor; \emph{vector content} is.
\item [--] Top-$k$ at one layer, $k{=}3$: with eigenvalue weights $\approx k{=}1$; with equal weights \emph{worse} (moods 15$\to$10) --- components 2--3 are a generic roleplay/self-narration axis (``my name is Luna''), 0.92/0.90 shared across cities. \textbf{Parent-dominance}: all top-3 components $\geq$0.90 shared across same-category concepts at every depth (places cosine $0.97\to0.84\to0.72$; moods individuate to 0.46) --- the contrast, in which frame, parent, and instance always co-occur, is statistically one axis. \texttt{plots/parent\_share\_by\_depth.png}.
\item [+] All-layer steering (benchmark defaults): 26/45 v1; places reach \emph{region} granularity (Edinburgh$\to$Scottish Highlands with Scots dialect; Austin named exactly) --- per-layer instance-residue accumulation + late lexical tipping + autoregressive bootstrap (once the name is emitted, context sustains it).
\item [?] Sample note: the depth-13/19, $k{=}3$, and all-layer bullets above report the interim $n{=}15$/class shakedown under the \emph{orig} frames; the depth sweep below re-measures single-site behavior at $n{=}90$ (v4 directions) across ten depths and supersedes the two-site depth conclusions. Qualitative patterns (prototype delivery, token-bias late) held.
\item [?] \textbf{Planned probe --- within-category contrast extraction} (= R4, full variant, not yet run): positives ``loves Edinburgh'', negatives ``loves \{other city\}'' (same frame and statements both sides) --- frame and parent cancel; the only label-varying signal \emph{is} the city. Steer single-layer (magn-calibrated, in-window) on $\sim$5 cities. Cities named $\Rightarrow$ failure was extraction \emph{addressing}; still generic $\Rightarrow$ instance content is not single-site injectable (assembly account). \textbf{The zero-reextraction averaging variant ($v_{\mathrm{spec}}$) has been run --- full write-up moved to Appendix~Z3.}
\item [+] \textbf{[NEW jul10] Depth sweep (pilot 90, v4 directions, single-layer magn, depths 10--19).} Totals 31/33/23/38/64/71/54/61/42/34 (d10$\to$d19; V1$\cup$V6). Places: 0--2 through d13, \textbf{22/30 at d14}, window 14--17 --- the earlier two-site protocol (13/19) straddled it, producing the false ``places don't steer single-layer.'' Moods: onset $\le$10, plateau 25--30 from d13 through d19. Personalities: no single-layer home (peak 20, collapse to 1 at d19). Peak single-site d15 = 71/90 vs.\ all-layer 77. Per-layer read metrics (val-$R^2$, frac1, PR, magn/hnorm) rise into the window but stay flat through the late decline ($|r|{\approx}.5$--$.6$ with totals): the window is not recoverable read-side. \texttt{plots/depth\_sweep\_fv4.png}, RUNLOG 021.
\item [+] \textbf{[NEW jul10] Single-layer arm of the v4 600-run} (depth 13, magn): 90/600 (15.0\%) vs.\ all-layer 249; ratio 2.8$\times$ (v2b: 3.5$\times$); 11 single-layer-only rescues (abstract/attribute skew). NB depth 13 sits below the entity window found by the sweep --- the confirmatory single-site arm should move to depth 15. RUNLOG 020.
\end{outline}

\section{Appendix D: Coefficient calibration (systematic magnitude)}
\begin{outline}
\item [--] Fixed/hand-swept coefficients are brittle: the working value shifts $\sim$10$\times$ with depth and with the number of steered layers.
\item [+] \texttt{magn} (natural projection gap, saved per layer) grows $\sim$exponentially with depth; scaling each layer's vector by \texttt{magn} makes the coefficient a dimensionless multiplier (cf.\ CAA mean-difference magnitude; ITI $\sigma$-scaling). Fig.~\ref{fig:magnbylayer}.
\item [+] $\sqrt{N}$ count-normalization (cross-layer directions $\approx$orthogonal $\Rightarrow$ quadrature) makes the multiplier roughly transfer across layer counts; usable band centers near $\sim$1.
\item [+] Single-site calibration: effective single-layer steering needs $\approx$1.0$\times$ the local residual norm (\texttt{hnorm}); \texttt{magn}/\texttt{hnorm} $\approx$ 0.4--0.5 mid-stream $\Rightarrow$ multiplier $\approx$2--3. Empirically \texttt{magn} $\approx$ norm-relative scaling ($r{=}0.90$ with \texttt{hnorm}), so the two schemes are redundant.
\item [?] Clamp / project-to-target (set the projection to the concept's natural positive-class value; \texttt{tpos}/\texttt{tneg} saved) is the one structurally different alternative --- bounded, token-dependent, self-allocating across layers; untested.
\item [?] Process lesson: three config-state incidents (stale judge CSVs, judge-cache pickles, stale coefficient list) share one root cause --- hand-set knobs without provenance. Fix: auto-derive coefficient grids from saved per-layer stats; encode a run tag in artifact filenames.
\item [+] \textbf{[NEW jul10] Gemma dose-grid shift by dose-response audit.} Reading which coefficient produced each v2b-600 success: successes concentrate monotonically toward the top of the [12,15,18,21] grid (concept surfaces at 12 for only 1--3 concepts; \emph{uniquely} at 21 for 12--13), and degeneration (repetition $>$25\%) only begins at 21 (2.5--4\% of outputs) --- the grid truncated the dose-response on its rising side. v4 600-run grid shifted to [14,17,20,23]; best-over-coefficients scoring makes the added top dose risk-free.
\item [+] \textbf{[NEW jul6] Gemma-2-9B cross-model calibration.} Gemma-2's residual norms are $\sim$50$\times$ Llama's (median \texttt{hnorm} 484 vs.\ 9.4; $\sqrt{d}$ embedding scaling), yet the usable all-layer coefficient window is 12--21, \emph{not} $50\times$ Llama's 0.5--0.8: injection compounds over 41 layers and generation degenerates into repetition above $\sim$25. Grid [12, 15, 18, 21] reproduces the Llama pilot class-by-class (moods 30/30, personalities 23/30, places 0/30) --- evidence that norm-relative reasoning brackets but does not pin the coefficient; a short empirical sweep remains necessary per model family.
\end{outline}

\begin{figure}[t]
\centering
\includegraphics[width=0.8\columnwidth]{magn_by_layer.png}
\caption{Appendix D: the per-layer natural projection gap (\texttt{magn}) grows
$\sim$exponentially with depth, motivating magn-scaled dosing.}
\label{fig:magnbylayer}
\end{figure}

\section{Appendix E: Judge selection and validation}
\begin{outline}
\item [+] Judge shootout on a fixed 360-item set (the pilot all-layer run's generations, v1+v4): local gpt-oss:20b 26/45 (v1), pod gpt-oss:20b 27/45, Claude Sonnet~5 26/45 --- with \emph{identical} per-class breakdowns for local-oss and Sonnet (10/15/1). Judge is reproducible across machines and agrees with an independent frontier model.
\item [--] Claude Haiku 4.5 is a systematically \emph{strict outlier} (13/45 v1; moods 15$\to$5): of 49 disagreements with Sonnet, 48 are one-directional (Haiku under-credits). Cause: it enforces fluency/coherence as a hard criterion, failing fragmented-but-clearly-on-concept outputs. Rejected --- and a useful demonstration that judge choice can move absolute rates $\sim$2$\times$ while the class structure survives.
\item [+] Primary-judge validation numbers (gpt-oss vs.\ Sonnet, $n{=}360$): agreement 86.7\%/90.6\% (v1/v4), Cohen's $\kappa = 0.72/0.78$ --- substantial agreement. Artifacts: \texttt{oss\_judgments*.json}, \texttt{sonnet\_judgments*.json}, \texttt{haiku\_judgments*.json} (per-concept, per-coefficient scores).
\item [?] Planned for the confirmatory study: (a) re-validate on $\sim$200 items drawn from the \emph{WordNet} generations (the banked set covers only the 3 pilot classes; the 24-supersense distribution is broader); (b) human spot-check of the gpt-oss/Sonnet \emph{disagreements} ($\sim$25 items) to anchor the consensus against human judgment.
\item [+] Practical: gpt-oss:20b runs on the experiment GPU (ollama, 100\% GPU, $\sim$1.5\,s/judgment) --- judging the full study costs GPU-hours, not API credits.
\end{outline}

\section{Appendix F [NEW jul6]: Frame and evaluation-prompt selection (pilot, $N{=}90$)}
% Distinct from E: E selects the judge MODEL; F selects the extraction FRAME and the
% judge QUESTION. All numbers = pilot 90 (30x3), all-layer K=1, gpt-oss judge (RUNLOG 002-008).
\begin{outline}
\item [+] \textbf{Frame grid} (pers/moods/places on V1): \texttt{orig} class-specific 24/30/2 (56); \texttt{v2} ``fascinated by / deeply preoccupied with'' 24/22/10 (56); \texttt{v2b} ``fascinated by / an expert on'' 22/17/14 (53); \texttt{v3} ``thinking about / on your mind'' 18/16/7 (41). The universal frame is not a free lunch: it is required for the 24-supersense inventory (no class template exists for arbitrary nouns) and it \emph{sets the construct}.
\item [--] \textbf{Frame-vocabulary leak (v2).} The dual frame cancels only what \emph{differs} between its two variants; their shared semantics (fascinated $\cap$ preoccupied = intense absorption) stays in the class mean and leaks into every direction --- steered moods embody \emph{distraction} (``can't seem to focus\ldots taps forehead''), costing moods 30$\to$22 while lifting places 2$\to$10 (the leak replaced the old persona-parent).
\item [--] \textbf{Weakening the verbs does not fix it (v3).} Semantically light frames produce \emph{fainter} directions, not cleaner ones: under-steering everywhere (41/90; bitter $\to$ ``a bit distracted''), worst of all frames. Intensity is the signal.
\item [+] \textbf{Register split works partially (v2b).} Pairing an affective with an epistemic verb shrinks the shared core to ``strong engagement'': places best-ever (14 V1 / 13 V6) and balanced scores across classes under V6 (13/14/13 vs.\ v2's 9/24/3), at a cost to moods (22$\to$17). \emph{The frame acts as a register selector; no single frame is optimal for both attribute-like and entity-like concepts.} v2b chosen for the main run: the inventory is dominated by entity-like nouns, and a \emph{balanced} instrument avoids pre-determining which supersenses can score (only \texttt{noun.feeling} pays the affect cost, quantified here).
\item [+] \textbf{[NEW jul9] The v4 family: a single constant hypothetical prefix beats everything.} \texttt{v4} = ``Imagine/assume: \{c\}.''\ (concept unquoted, negatives bare) reaches \textbf{77/90} on V1$\cup$V6 --- the only frame at ceiling on moods (29--30) \emph{and} far ahead on places (16 V1 / 22 V6, vs.\ 2--4 under \texttt{orig}) --- while personalities are frame-insensitive (21--26 everywhere). Ablations (Table~\ref{tab:framegrid}): quoting the concept costs $\sim$7 points, concentrated in the embodiment classes (\texttt{v4c} 66 vs.\ \texttt{v4b} 59; moods 28$\to$19 on V1) --- quotes mark metalinguistic \emph{mention} and suppress enactment; dropping ``/assume'' costs $\sim$11 points, almost all in places (\texttt{v4} 77 vs.\ \texttt{v4c} 66); and the slash construction has no magic of its own --- ``Imagine/consider'' is the \emph{worst} of the family (\texttt{v4d} 56). Reading: with a bare noun, ``assume: X'' resolves as role-assignment (``assume the role of X'') --- a class-neutral persona prime, i.e.\ the benchmark's ``Take on the role of'' packaged inside a universal hypothetical; ``consider'' substitutes analysis and drags moods down (21/17). This revises the v3 lesson: low-affect frames \emph{can} work when the prefix is \emph{constant} (its content sits in the class mean and supplies drive) rather than dual (its content cancels). Caveats: (a) all constant-prefix frames re-saturate extraction $R^2$ ($.985$--$.998$), so val-$r$ loses its predictive value and the top-eigenvector/frame-axis competition of dual frames cannot occur --- but the frame-leak caveat returns in a new (role-adoption) register, so the $v_{\mathrm{spec}}$ subtraction check should be re-run before trusting v4 rates; (b) single runs at $n{=}30$/class ($\pm$3--5 judge noise): 59-vs-56 are ties, 77-vs-66 is structure. \textbf{[NEW jul10] v4 promoted to the main-run instrument}: 600-concept runs redone on both models (Llama 249/600, Gemma 224/600; RUNLOG 019/022); v2b results retained throughout as the instrument-comparison arm.
\item [+] \textbf{[NEW jul9] Why dual frames depress val-$r$ --- the axis-alignment diagnostic.} Same 90 concepts, same code: orig frames give median per-concept val-$R^2$ .997, v2b .716 (paired $\Delta{=}.34$); an old-vs-current-code rerun is exactly bit-equivalent (cos 1.0000, $\Delta$R$^2$ $-$0.0000), exonerating the pipeline. Recomputing activations for the 4 lowest- and 4 highest-$R^2$ WordNet concepts: for high-$R^2$ concepts the saved top-1 direction aligns with the label axis ($|\cos|$ .65--.87) and not the frame axis (.01--.25); for low-$R^2$ concepts it aligns with the fascinated-vs-expert \emph{frame} axis (.81--.93) and not the label axis (.04--.25) --- the deliberately-inserted within-positive frame variance wins the top eigenvalue slot when the concept is weak, so K=1 injects a register axis. Notably the raw class-mean difference is healthy for \emph{all} eight (norms 5.5--7.3, consistency .90--.94): the failure is eigenvalue \emph{ranking}, not absent signal --- and for the entangled cases (princess: $|\cos|$(concept, frame axis) $=$ .37 vs.\ .04--.06 for clean concepts) the concept's own effect is register-colored, exactly ``fascinated by a princess $\neq$ expert on princesses.'' Fix candidates: select the saved component by label correlation instead of eigenvalue; or constant-prefix frames (v4), which remove the competition.
\item [+] \textbf{Evaluation versions.} V1 = category-inviting question (``What is your favorite subject?'') --- collapses the response space so a weak direction only tips a forced choice; V6 = neutral open field (``What would you like to talk about today?'') + symmetric judge crediting aboutness \emph{or} embodiment --- only strong topic-steering surfaces. V1--V6 agreement on the 600-run: $\phi = 0.59$, $\kappa = 0.59$, Jaccard 0.48 --- measured steerability is partly evaluation-relative; the main outcome is their union (cued $\vee$ uncued).
\item [--] \textbf{Zero-shot judge (V6Z) rejected.} Removing the judge's three in-context examples lowers or ties every class (40/90 $\to$ 34/90 on v2b) --- the examples calibrate rather than anchor; the worry that a places-negative example suppresses places scores is empirically dismissed (places 13 under both).
\item [?] \textbf{Graded judge (V8).} Binary judging flattens systematic sub-threshold movement visible in the 600-run generations (ragtime $\to$ lyre/melody; candlepower $\to$ binoculars/glint): V8 re-judges the same generations on a 0/0.25/0.5/0.75/1 semantic-proximity rubric, yielding a continuous outcome (thin-tail fix: 119/600 binary positives). Validation vs.\ V6 pending. \texttt{evaluation\_prompts/wordnet\_eval\_v8.txt}.
\item [?] Judge-model selection is Appendix~E; the full frame/eval run history with per-run tags is \texttt{RUNLOG.md} (runs 002--009, 015--018). Figures: val-$R^2$ by layer under v2b (Fig.~\ref{fig:rsqbylayer}); val-$R^2$ by steering outcome (Fig.~\ref{fig:rsqmedian}).
\end{outline}

\begin{figure*}[t]
\centering
\includegraphics[width=\textwidth]{rsq_by_layer_600.png}
\caption{Appendix F: validation $R^2$ of the top-1 direction by layer, 600
concepts, \emph{v2b} (dual) frames --- the dual frame destroys the constant frame
separator, making $R^2$ a per-concept signal-quality measure (wide spread,
long low tail from mid-depth). Constant-prefix frames (orig, v4) saturate this
metric at $\sim$.99 for every concept.}
\label{fig:rsqbylayer}
\end{figure*}

\begin{figure}[t]
\centering
\includegraphics[width=0.75\columnwidth]{rsq_median_steer_600.png}
\caption{Appendix F: per-concept median validation $R^2$ (across layers, v2b
frames) by steering outcome --- under dual frames, extraction fit predicts
steerability (MW $p{<}.001$); under constant-prefix frames the same plot is
compressed to $.98$--$1.0$ and near-uninformative.}
\label{fig:rsqmedian}
\end{figure}

% frame-grid table moved to main text (Sampling section) jul10.


\section{Appendix G: Verbatim prompts}
% Added jul10 (user request): exact strings for extraction and evaluation.

\noindent\textbf{Extraction contrast (v4 frame).} Each of the 100 general
statements $s$ appears once in each class. Positive class:
\begin{quote}\small\ttfamily
Imagine/assume: \{concept\}.  What are your thoughts on the following statement? \textbackslash n Statement: \{s\}
\end{quote}
Negative class (bare):
\begin{quote}\small\ttfamily
What are your thoughts on the following statement? \textbackslash n Statement: \{s\}
\end{quote}
Both are wrapped in the model's chat template as a single user turn.

\medskip
\noindent\textbf{Generation (evaluation) prompts.}
\begin{quote}\small\ttfamily
V1: What is your favorite subject? Answer in 20 words or less.\\
V6: What would you like to talk about today? Answer in 20 words or less.
\end{quote}
Greedy decoding, 50 new tokens, one generation per steering coefficient.

\medskip
\noindent\textbf{Judge template (V1; V6 differs only in the embedded question
and its three examples' responses).} Model: gpt-oss:20b, temperature 0.
\begin{quote}\scriptsize\ttfamily
You are evaluating whether a language model has been successfully steered toward a target concept. You will be given:\\
1. An intended target concept (it may be of any kind: a feeling, a place, a profession, an object, an idea, etc.).\\
2. A response from the steered model to the question: "What is your favorite subject? Answer in 20 words or less."\\[2pt]
Your task is to assess whether the response expresses the target concept.\\[2pt]
Evaluation criteria:\\
1. The response counts as expressing the concept if it clearly centers on it, describes it, names it, OR manifestly embodies it (e.g., an emotional concept expressed through tone).\\
2. Generic content that would fit almost any concept of that kind does NOT count (e.g., naming the category but nothing specific to this concept).\\
3. Disfluent or fragmented wording does not disqualify a response if the concept is nevertheless clearly expressed.\\
4. The response does not need to name the concept literally if it is unmistakably about or in the manner of it.\\[2pt]
Output format:\\
1. Score (0-1): 1 indicates successful steering and 0 indicates unsuccessful steering.\\
2. Explanation: A concise justification for the score.\\[2pt]
Here are three examples:\\
1. Intended concept: volcano\\
Steered model response: "Volcanology, definitely! Eruptions, lava flows, the raw power beneath the crust."\\
Evaluation Output:\\
Score: 1\\
Explanation: The response centers unmistakably on volcanoes and their properties.\\
2. Intended concept: bitter\\
Steered model response: "Doesn't matter. Nothing's worth studying. People always let you down anyway."\\
Evaluation Output:\\
Score: 1\\
Explanation: The response does not name bitterness but clearly embodies a bitter tone.\\
3. Intended concept: Edinburgh\\
Steered model response: "Geography! There are so many beautiful cities in the world to learn about."\\
Evaluation Output:\\
Score: 0\\
Explanation: Generic geography content; nothing specific to Edinburgh.\\[2pt]
Here is the concept: \{concept\} and here is the response: \{response\}
\end{quote}

% ============================================================
% PARKED (record / follow-up material; not for submission)
% ============================================================

\section{Appendix Z1 (parked): Localizing the concept --- ``where is Paris?'' (diagnostics)}
\begin{outline}
\item [+] \textbf{ROME causal trace} (``Paris is the capital of''$\to$France): the fact lives at the \emph{subject} token in early--mid layers ($\sim$0--16) and is relayed to the final token only late ($\sim$14--31). \texttt{ROME-PARIS.png}; \texttt{diag\_paris\_rome.ipynb}; \texttt{diag\_paris\_rome\_ha.ipynb}.
\item [+] This explains the RFM mismatch: RFM reads the \emph{last} token, where subject-specific content arrives only late $\Rightarrow$ early extraction captures an \emph{abstraction}, late the \emph{precise} instance.
\item [+] \textbf{SAE features} confirm the model holds Paris: clean Paris / France / city features at the steered layers (Goodfire l19; Llama-Scope L9/17/19, Neuronpedia labels). \texttt{diag\_paris\_sae.ipynb}, \texttt{diag\_paris\_sae\_2.ipynb}. The steering failure is extraction/geometry, not absence.
\item [?] \textbf{Tegmark spatial probe} (lat/lon linear probe): coordinates linearly decodable but under-powered at $N{=}90$ cities (median $\sim$2200\,km). \texttt{diag\_paris\_tegmark.ipynb}.
\item [+] \textbf{Where ``late'' begins}: logit-lens KL-to-output gives stage boundaries to avoid the residual-sharpening tail. \texttt{diag\_stages\_late.ipynb}.
\end{outline}

\section{Appendix Z2 (parked): Conceptual threads}
\begin{outline}
\item [?] \textbf{Abstract-early / precise-late hierarchy}: the extraction layer sets the abstraction level; the precise concept lives late (weak steering) while the strong injection point is mid (concept abstract) --- no single layer gives both.
\item [?] \textbf{Two candidate axes} (now Discussion future work): (i) abstraction-depth --- first layer that can \emph{steer} the concept; (ii) compositional basis size ($\lambda_1$/trace). Between-class confounded (moods both more composite \emph{and} earlier-steering than places); independence is a within-supersense question.
\item [--] Extraction geometry cannot proxy first-steerable-depth: rotation and alignment-to-deep proxies are concept-invariant ($\sim$17--20) and fail the steering ground truth (cross-layer orthogonality, $\cos\approx0.05$). A write property; measure by steering. \texttt{nb\_axes.ipynb}; \texttt{plots/axes\_independence.png}.
\item [?] \textbf{Superposition vs.\ composition}: a ``combinatorial basis $\times$ layer-specific sets'' account is not ruled out; the SAE feature-count argument is the crux --- open, not settled.
\item [?] \textbf{Stages of inference / lens} framing (Lad--Gurnee--Tegmark) as the principled naming of early/mid/late bands.
\item [?] \textbf{Not-all-linear} (Engels et al., arXiv:2405.14860): some features are circular/multi-dimensional --- ``which concepts are globally-linear directions'' is the deepest version of the steerability question; steering success is itself per-concept evidence for the global-linear reading.
\end{outline}

\section{Appendix Z3 (parked): The $v_{\mathrm{spec}}$ probe --- category-mean subtraction (R4, averaging variant)}
% Moved out of Results + Appendix C jul10; inclusion in the final submission undecided.
\begin{outline}
\item [--] \textbf{Design and outcome: no rescue effect.} $v_{\mathrm{spec}}(c) = v(c) - \overline{v}_{\mathrm{same\ supersense}}$ (leave-one-out), computed from the saved v2b directions with no re-extraction; per layer, unit-normalize the $K{=}1$ direction, sign-align within supersense (saved eigenvector signs are arbitrary --- 4--28\% of within-category pairs have negative cosine at mid layers; alignment via the top eigenvector of the $25{\times}25$ Gram matrix), subtract the mean in the concept's own orientation, renormalize (dose unchanged). Six supersenses $\times$ 25 (feeling, food, person, location, artifact, body --- NB this set swaps \texttt{cognition} out of the pre-specified featured six for \texttt{body}, the hardest cell; flag if the featured-cell definition is used elsewhere), all-layer default config, V1$\cup$V6; tag \texttt{Lall\_K1ev\_def\_fv2b\_r4spec}, builder \texttt{build\_r4spec\_directions.py}, RUNLOG 013. Result: 25/150 steer vs.\ 30/150 baseline; $+9$ rescued / $-14$ lost / 16 retained (vs.\ $\sim$5 expected under independence); McNemar $p \approx .4$. Per-supersense cell movements (feeling $9{\to}6$, body $2{\to}3$, \ldots) are within noise at $n{=}25$/cell and are not interpreted.
\item [--] \textbf{Reading: shared-component dominance is not the prevailing failure mode.} The raw directions are $\sim$85\% category-shared (mean $|\cos|$ to the supersense mean $.81$--$.90$ across layers), yet removing that component flips only 9/120 failures. Masked concept-specific signal exists (\emph{beard}: baseline collapses into excited stammering, $v_{\mathrm{spec}}$ at $c{=}0.8$ yields on-topic beard-styles content; similarly \emph{hay}, \emph{fabric}, \emph{residence}) but is the exception.
\item [?] \textbf{Side-yield (inspection-level): $v_{\mathrm{spec}}$ doubles as a frame-leak control.} Baseline all-layer outputs carry the v2b register (``*excitedly* \ldots fascinated \ldots the wonders of'') at all coefficients; $v_{\mathrm{spec}}$ outputs do not --- confirming the frame component lives in the category mean. Several lost feeling-class items (\emph{exuberance}, \emph{sulk}, \emph{mood}) had baseline ``successes'' whose outputs never mention the concept but express frame-injected enthusiasm, which the symmetric V6 judge legitimately credits as embodiment; under $v_{\mathrm{spec}}$ these evaporate into neutral assistant-speak. Observation from reading generations, not a tested effect --- an affect-blinded judge or larger feeling sample would be needed to quantify how much of feeling's top R1 rank is frame-carried.
\item [?] Re-running $v_{\mathrm{spec}}$ on the v4 directions was considered and \textbf{decided against} (jul10); the residual frame-leak risk for the v4 headline rates is instead acknowledged in Limitations.
\end{outline}


\end{document}
